\documentclass[aps,preprint]{revtex4}%
\usepackage{amsfonts}
\usepackage{amsmath}
\usepackage{amssymb}
\usepackage{graphicx}%
\setcounter{MaxMatrixCols}{30}
%TCIDATA{OutputFilter=latex2.dll}
%TCIDATA{Version=5.50.0.2953}
%TCIDATA{CSTFile=revtex4.cst}
%TCIDATA{Created=Saturday, July 17, 2004 08:28:44}
%TCIDATA{LastRevised=Tuesday, July 24, 2007 10:43:53}
%TCIDATA{<META NAME="GraphicsSave" CONTENT="32">}
%TCIDATA{<META NAME="SaveForMode" CONTENT="1">}
%TCIDATA{BibliographyScheme=Manual}
%TCIDATA{<META NAME="DocumentShell" CONTENT="Articles\SW\REVTeX 4">}
%BeginMSIPreambleData
\providecommand{\U}[1]{\protect\rule{.1in}{.1in}}
%EndMSIPreambleData
\newtheorem{theorem}{Theorem}
\newtheorem{acknowledgement}[theorem]{Acknowledgement}
\newtheorem{algorithm}[theorem]{Algorithm}
\newtheorem{axiom}[theorem]{Axiom}
\newtheorem{claim}[theorem]{Claim}
\newtheorem{conclusion}[theorem]{Conclusion}
\newtheorem{condition}[theorem]{Condition}
\newtheorem{conjecture}[theorem]{Conjecture}
\newtheorem{corollary}[theorem]{Corollary}
\newtheorem{criterion}[theorem]{Criterion}
\newtheorem{definition}[theorem]{Definition}
\newtheorem{example}[theorem]{Example}
\newtheorem{exercise}[theorem]{Exercise}
\newtheorem{lemma}[theorem]{Lemma}
\newtheorem{notation}[theorem]{Notation}
\newtheorem{problem}[theorem]{Problem}
\newtheorem{proposition}[theorem]{Proposition}
\newtheorem{remark}[theorem]{Remark}
\newtheorem{solution}[theorem]{Solution}
\newtheorem{summary}[theorem]{Summary}
\newenvironment{proof}[1][Proof]{\noindent\textbf{#1.} }{\ \rule{0.5em}{0.5em}}
\begin{document}
\preprint{HEP/123-qed}
\title[Short title for running header]{Long title that appears on the title page}
\author{First Author}
\affiliation{My Institution}
\author{Second Author}
\affiliation{My Institution}
\author{Third Author}
\affiliation{Other Institution}
\keywords{one two three}
\pacs{PACS number}

\begin{abstract}
Shell document for REV\TeX{} 4.

\end{abstract}
\volumeyear{year}
\volumenumber{number}
\issuenumber{number}
\eid{identifier}
\date[Date text]{date}
\received[Received text]{date}

\revised[Revised text]{date}

\accepted[Accepted text]{date}

\published[Published text]{date}

\startpage{101}
\endpage{102}
\maketitle
\tableofcontents


\subsection{Properties of an Optimized AGP\ Fock Operator\label{AGP_Fock}}

The generalized Fock operator associated with an optimized AGP\ state has
further significant properties, in addition to being self adjoint, originating from

\begin{enumerate}
\item the stationarity wrt occupation number variations shown in eq. $\left(
\ref{occnumstat}\right)  $ as
\begin{equation}
\frac{1}{S_{N}\left(  g^{N}\right)  }\left\langle \left.  g^{N}\right\vert
Ha_{p}^{\dagger}a_{p}g^{N}\right\rangle =\mathcal{E}\left(  g^{N}\right)
N_{p};\quad1\leq p\leq2s,
\end{equation}
and

\item the expression for the generalized Fock operator matrix elements
\begin{equation}
F\left(  \mathbb{D}^{2},\mathbf{\varphi}\right)  _{qp}=\frac{1}%
{\operatorname{Tr}\left\{  D^{2N}\right\}  }\operatorname{Tr}\left\{
a_{p}^{\dagger}\left[  a_{q},H\right]  D^{2N}\right\}
\end{equation}
obtained in Appendix \ref{SQ_Fock}, which are valid for any state described by
the positive operator $D^{2N}$. For AGP\ states this becomes
\begin{equation}
F\left(  \mathbf{n},\mathbf{\psi}\right)  _{qp}=\frac{1}{S_{N}\left(
g^{N}\right)  }\left\langle \left.  g^{N}\right\vert a_{p}^{\dagger}%
a_{q}H-a_{p}^{\dagger}Ha_{q}g^{N}\right\rangle .
\end{equation}

\end{enumerate}

Combining items $\left(  1\right)  $ and $\left(  2\right)  $ together shows
that the diagonal elements of the optimized AGP\ Fock operator are given by
\begin{equation}
F\left(  \mathbf{n},\mathbf{\psi}\right)  _{pp}=\frac{1}{S_{N}\left(
g^{N}\right)  }\left\langle \left.  g^{N}\right\vert a_{p}^{\dagger}%
a_{p}H-a_{p}^{\dagger}Ha_{p}g^{N}\right\rangle =\mathcal{E}\left(
g^{N}\right)  N_{p}-\frac{1}{S_{N}\left(  g^{N}\right)  }\left\langle \left.
a_{p}g^{N}\right\vert Ha_{p}g^{N}\right\rangle ,
\end{equation}
which can be developed to give
\begin{equation}
F\left(  \mathbf{n},\mathbf{\psi}\right)  _{pp}=N_{p}\left\{  \mathcal{E}%
\left(  g^{N}\right)  -\frac{\left\langle \left.  a_{p}g^{N}\right\vert
Ha_{p}g^{N}\right\rangle }{\left\langle \left.  a_{p}g^{N}\right\vert
a_{p}g^{N}\right\rangle }\right\}  =N_{p}\left\{  \mathcal{E}\left(
g^{N}\right)  -\mathcal{E}_{2N-1}\left(  \psi_{p}\right)  \right\}  ,
\end{equation}
where $\mathcal{E}_{2N-1}\left(  \psi_{p}\right)  $ is the energy of the
normalized $2N-1$ electron state produced from the state $\frac{1}%
{S_{N}\left(  g^{N}\right)  }\left\vert g^{N}\right\rangle $ by removing an
electron described by the CGSO\ $\left\vert \psi_{p}\right\rangle $. One can
observe from the form of $F\left(  \mathbf{n},\mathbf{\psi}\right)  _{pp}$
that
\begin{equation}
F\left(  \mathbf{n},\mathbf{\psi}\right)  _{pp}=F\left(  \mathbf{n}%
,\mathbf{\psi}\right)  _{p+sp+s}\Rightarrow I\left(  \psi_{p}\right)
=I\left(  \psi_{p+s}\right)  ;\quad1\leq p\leq s.
\end{equation}
In Appendix \ref{Fock&Energy} we show that the total energy can be expressed
in terms of the Fock operator $F\left(  \mathbf{n},\mathbf{\psi}\right)
\ $as
\begin{equation}
\mathcal{E}\left(  \mathbf{n},\mathbf{\psi}\right)  =\frac{1}{2}\left\{
\overset{}{\operatorname{Tr}\left\{  F\left(  \mathbf{n},\mathbf{\psi}\right)
\right\}  +\mathcal{E}_{1}\left(  \mathbf{n},\mathbf{\psi}\right)  }\right\}
,
\end{equation}
where the expectation value $\mathcal{E}_{1}\left(  \mathbf{n},\mathbf{\psi
}\right)  $ of the one electron part of the Hamiltonian is given by
\begin{equation}
\mathcal{E}_{1}\left(  \mathbf{n},\mathbf{\psi}\right)  =\sum_{1\leq p\leq
2s}N_{p}\left(  \left.  \psi_{p}\right\vert h_{1}\psi_{p}\right)  ,
\end{equation}
as the FORDM $\mathbb{D}^{1}\left(  \mathbf{n},\mathbf{\psi}\right)  $ is
diagonal in the CGSO's $\left\{  \left\vert \psi_{p}\right\rangle \right\}  .$
The energy $\mathcal{E}\left(  \mathbf{n},\mathbf{\psi}\right)  $ can be
expressed as
\begin{align}
\mathcal{E}\left(  \mathbf{n},\mathbf{\psi}\right)   &  =\frac{1}{2}%
\sum_{1\leq p\leq2s}\left\{  N_{p}I\left(  \psi_{p}\right)  +N_{p}\left(
\left.  \psi_{p}\right\vert h_{1}\psi_{p}\right)  \right\} \\
&  =\frac{1}{2}\sum_{1\leq p\leq2s}\left\{  N_{p}I\left(  \psi_{p}\right)
+N_{p}\left(  \left.  \psi_{p}\right\vert h_{1}\psi_{p}\right)  \right\}
\equiv\frac{1}{2}%
%TCIMACRO{\tsum \limits_{1\leq l\leq2s}}%
%BeginExpansion
{\textstyle\sum\limits_{1\leq l\leq2s}}
%EndExpansion
N_{l}\varsigma\left(  \psi_{l}\right)  ,
\end{align}
where we have introduced the new quantity,
\begin{equation}
\varsigma\left(  \psi_{l}\right)  =I\left(  \psi_{l}\right)  +\left(  \left.
\psi_{l}\right\vert h_{1}\psi_{l}\right)  ,
\end{equation}
which is the sum of the ionization potential, the kinetic energy, the electron
- nuclear potential energy of attraction and potential energy of interaction
due to any other external forces of an electron described by a
GSO\ $\left\vert \psi_{l}\right\rangle $.

Although the AGP\ Fock operator is self adjoint its matrix, which is
hermitian, in the CGSO basis is in general non diagonal. If the canonical
coefficients of the optimized AGP\ are non degenerate then the generalized
Brillouin conditions eq. $\left(  \ref{GBC}\right)  $ hold and
\begin{equation}
F\left(  \mathbf{n},\mathbf{\psi}\right)  _{qp}=\frac{1}{S_{N}\left(
g^{N}\right)  }\left\{  \left\langle \left.  g^{N}\right\vert a_{p}^{\dagger
}a_{p}Hg^{N}\right\rangle \delta_{qp}-\left\langle \left.  a_{p}%
g^{N}\right\vert Ha_{q}g^{N}\right\rangle \right\}  .
\end{equation}
The matrix $\left\{  F\left(  \mathbf{n},\mathbf{\psi}\right)  _{qp}\right\}
$ can be diagonalized by a unitary transformation $u$ of the CGSO basis to
produce
\begin{equation}
\tilde{F}\left(  \mathbf{n},\mathbf{\psi}\right)  _{qp}=\frac{1}{S_{N}\left(
g^{N}\right)  }\left\{  \left\langle \left.  g^{N}\right\vert b_{p}^{\dagger
}b_{p}Hg^{N}\right\rangle -\left\langle \left.  b_{p}g^{N}\right\vert
Hb_{p}g^{N}\right\rangle \right\}  \delta_{qp},
\end{equation}
where
\begin{equation}
b_{p}^{\dagger}=ua_{p}^{\dagger}u^{\dagger}=\sum_{1\leq j\leq2s}a_{j}%
^{\dagger}u_{jp}%
\end{equation}
and
\begin{equation}
\left\vert \varsigma_{p}\right\rangle =\sum_{1\leq j\leq2s}\left\vert \psi
_{j}\right\rangle u_{jp}.
\end{equation}
The eigenvalues
\begin{align}
\varepsilon_{p}  &  =\tilde{F}\left(  \mathbf{n},\mathbf{\psi}\right)
_{pp}=\frac{1}{S_{N}\left(  g^{N}\right)  }\left\{  \left\langle \left.
g^{N}\right\vert b_{p}^{\dagger}b_{p}Hg^{N}\right\rangle -\left\langle \left.
b_{p}g^{N}\right\vert Hb_{p}g^{N}\right\rangle \right\} \nonumber\\
&  =\frac{1}{S_{N}\left(  g^{N}\right)  }\mathcal{N}_{p}\left\{
\mathcal{E}\left(  g^{N}\right)  -\frac{\left\langle \left.  b_{p}%
g^{N}\right\vert Hb_{p}g^{N}\right\rangle }{\left\langle \left.  b_{p}%
g^{N}\right\vert b_{p}g^{N}\right\rangle }\right\}  ,
\end{align}
where
\begin{equation}
\mathcal{N}_{p}=\left\langle \left.  \varsigma_{p}\right\vert D^{1}\left(
\mathbf{n},\mathbf{\psi}\right)  \varsigma_{p}\right\rangle .
\end{equation}
Hence
\begin{equation}
\varepsilon_{p}=\frac{\mathcal{N}_{p}I\left(  \varsigma_{p}\right)  }%
{S_{N}\left(  g^{N}\right)  },
\end{equation}
where $\left\{  I\left(  \varsigma_{p}\right)  \right\}  $ are the ionization
potentials associated with the GSO's $\left\{  \left\vert \varsigma
_{p}\right\rangle \right\}  .$ The optimized geminal\ \emph{is not} in general
in canonical form wrt the basis $\left\{  \left\vert \varsigma_{p}%
\right\rangle \right\}  $ , nor is $D^{1}\left(  \mathbf{n},\mathbf{\psi
}\right)  $ diagonal so $\mathcal{N}_{p}$ are \emph{not} occupation
numbers.Noting that
\begin{align}
&  U\left(  \mathbf{\alpha}\right)  ^{\dagger}a_{i}^{\dagger}U\left(
\mathbf{\alpha}\right)  =\sum a_{j}^{\dagger}U\left(  \mathbf{\alpha}\right)
_{ji}^{\ast}=\sum_{1\leq j\leq2s}U\left(  \mathbf{\alpha}\right)  _{ij}%
a_{j}^{\dagger}\\
&  U\left(  \mathbf{\alpha}\right)  ^{\dagger}a_{i}^{\dagger}U\left(
\mathbf{\alpha}\right)  \left\vert \phi\right\rangle =U\left(  \mathbf{\alpha
}\right)  ^{\dagger}a_{i}^{\dagger}\left\vert \phi\right\rangle =U\left(
\mathbf{\alpha}\right)  ^{\dagger}\left\vert \varphi_{i}\right\rangle
=\sum_{j}\left\vert \varphi_{j}\right\rangle U\left(  \mathbf{\alpha}\right)
_{ji}^{\dagger}=\sum_{j}\left\vert \varphi_{j}\right\rangle U\left(
\mathbf{\alpha}\right)  _{ij}^{\ast}\\
&  U\left(  \mathbf{\alpha}\right)  ^{\dagger}a_{k}U\left(  \mathbf{\alpha
}\right)  =\sum a_{l}U\left(  \mathbf{\alpha}\right)  _{kl}^{\ast}=\sum_{1\leq
l\leq2s}a_{l}U\left(  \mathbf{\alpha}\right)  _{lk}\\
&  \sum_{1\leq i,k\leq2s}U\left(  \mathbf{\alpha}\right)  _{ij}h_{1ik}U\left(
\mathbf{\alpha}\right)  _{lk}a_{j}^{\dagger}a_{l}=\sum_{1\leq i,j,k,l\leq
2s}U\left(  \mathbf{\alpha}\right)  _{ji}^{\ast}\left\langle \left.
\varphi_{i}\right\vert h_{1}\varphi_{k}\right\rangle U\left(  \mathbf{\alpha
}\right)  _{lk}a_{j}^{\dagger}a_{l}\\
&  =\sum_{1\leq i,j,k,l\leq2s}\left\langle \left.  U\left(  \mathbf{\alpha
}\right)  _{ji}\varphi_{i}\right\vert h_{1}U\left(  \mathbf{\alpha}\right)
_{lk}\varphi_{k}\right\rangle a_{j}^{\dagger}a_{l}=\sum_{1\leq i,j,k,l\leq
2s}\left\langle \left.  \varphi_{j}\left(  \mathbf{\alpha}\right)  \right\vert
h_{1}\varphi_{l}\left(  \mathbf{\alpha}\right)  \right\rangle a_{j}^{\dagger
}a_{l}%
\end{align}
%

\begin{align}
\mathcal{E}\left(  \mathbf{\psi}\right)   &  =\operatorname{Tr}\left\{
\left\{  \frac{1}{\left(  2N-1\right)  }\sum_{1\leq i,j,k\leq2s}h_{1ij}%
a_{i}^{\dagger}a_{k}^{\dagger}a_{k}a_{j}+\frac{1}{4}\sum_{1\leq i,j,k,l\leq
2s}h_{2ijkl}a_{i}^{\dagger}a_{j}^{\dagger}a_{k}a_{l}\right\}  D^{2N}\right\}
\\
&  =\operatorname{Tr}\left\{  \left\{  \frac{1}{\left(  2N-1\right)  }%
\sum_{1\leq i,j\leq2s}h_{1ij}a_{i}^{\dagger}a_{k}^{\dagger}a_{k}a_{j}+\frac
{1}{4}\sum_{1\leq i,j,k,l\leq2s}h_{2ijkl}a_{i}^{\dagger}a_{j}^{\dagger}%
a_{k}a_{l}\right\}  D^{2N}\right\}
\end{align}
One can observe from the form of $F\left(  \mathbf{n},\mathbf{\psi}\right)
_{pp}$ that
\begin{equation}
F\left(  \mathbf{n},\mathbf{\psi}\right)  _{pp}=F\left(  \mathbf{n}%
,\mathbf{\psi}\right)  _{p+sp+s}\Rightarrow I\left(  \psi_{p}\right)
=I\left(  \psi_{p+s}\right)  ;\quad1\leq p\leq s.
\end{equation}%
\begin{equation}
D_{klij}^{2}=\operatorname{Tr}\left\{  a_{i}^{\dagger}a_{j}^{\dagger}%
a_{k}a_{l}D^{2N}\right\}
\end{equation}
Fix $\left\{  D_{klij}^{2}\right\}  $ and consider
\begin{align}
\mathcal{E}\left(  \mathbf{\psi}\right)   &  =\frac{1}{4}\sum_{1\leq
i,j,k,l\leq2s}\left\{  \overset{}{\frac{1}{\left(  2N-1\right)  }\left\{
h_{1ik}\delta_{jl}+h_{1jl}\delta_{ik}-h_{1il}\delta_{jk}-h_{1jk}\delta
_{il}\right\}  +h_{2ijkl}}\right\}  D_{klij}^{2}\\
&  =\frac{1}{4}\sum_{1\leq i,j,k,l\leq2s}\left\{  \frac{2}{\left(
2N-1\right)  }\overset{}{\left\{  \left\langle \left.  \psi_{i}\right\vert
h_{1}\psi_{k}\right\rangle \delta_{jl}-\left\langle \left.  \psi
_{i}\right\vert h_{1}\psi_{l}\right\rangle \delta_{jk}\right\}  }\right. \\
&  \qquad\qquad\qquad\qquad\qquad\left.  \overset{}{}+\left\{  \left(  \left.
\psi_{i}\psi_{j}\right\vert h_{2}\psi_{k}\psi_{l}\right)  -\left(  \left.
\psi_{i}\psi_{j}\right\vert h_{2}\psi_{l}\psi_{k}\right)  \right\}  \right\}
D_{klij}^{2}%
\end{align}%
\begin{align}
\mathcal{E}\left(  \mathbf{\psi}\right)   &  =\frac{1}{2}\sum_{1\leq
i,j,k,l\leq2s}\left\{  \frac{1}{\left(  2N-1\right)  }\overset{}{\left\{
\left\langle \left.  \psi_{i}\right\vert h_{1}\psi_{k}\right\rangle
\delta_{jl}-\left\langle \left.  \psi_{i}\right\vert h_{1}\psi_{l}%
\right\rangle \delta_{jk}\right\}  }\right. \\
&  \qquad\qquad\qquad\qquad\qquad+\left.  \left\{  \left(  \left.  \psi
_{i}\psi_{j}\right\vert h_{2}\psi_{k}\psi_{l}\right)  -\left(  \left.
\psi_{i}\psi_{j}\right\vert h_{2}\psi_{l}\psi_{k}\right)  \right\}  \right\}
D_{klij}^{2}%
\end{align}
The reduced Hamiltonian $K$ is defined as
\begin{align}
K  &  =\frac{1}{2}\sum_{1\leq i,j,k,l\leq2s}\left\{  \frac{1}{\left(
2N-1\right)  }\overset{}{\left\{  \left\langle \left.  \psi_{i}\right\vert
h_{1}\psi_{k}\right\rangle \delta_{jl}-\left\langle \left.  \psi
_{i}\right\vert h_{1}\psi_{l}\right\rangle \delta_{jk}\right\}  }\right. \\
&  \qquad\qquad+\left.  \overset{}{\left\{  \left(  \left.  \psi_{i}\psi
_{j}\right\vert h_{2}\psi_{k}\psi_{l}\right)  -\left(  \left.  \psi_{i}%
\psi_{j}\right\vert h_{2}\psi_{l}\psi_{k}\right)  \right\}  }\right\}
\left\vert \psi_{i}\psi_{j}\right\rangle \left\langle \psi_{l}\psi
_{k}\right\vert ,
\end{align}
giving
\begin{equation}
\mathcal{E}\left(  \mathbf{\psi}\right)  =\operatorname{Tr}\left\{  K\left(
\mathbf{\psi}\right)  D^{2}\right\}  =\frac{1}{4}\sum_{1\leq i,j,k,l\leq
2s}K_{ijkl}\left(  \mathbf{\psi}\right)  D_{klij}^{2}.
\end{equation}
The functional derivatives are
\begin{equation}
\frac{\delta}{\delta\psi_{\mu}^{\ast}}\left\langle \left.  \psi_{a}\right\vert
h_{1}\psi_{b}\right\rangle =\left\{  \sum_{1\leq m,n\leq2s}h_{1mb}\delta
_{an}\left\vert \psi_{m}\right\rangle \left\langle \psi_{n}\right\vert
\right\}  \left\vert \psi_{\mu}\right\rangle ,
\end{equation}
and%
\begin{align}
&  \frac{\delta}{\delta\psi_{\mu}^{\ast}}\left(  \left.  \psi_{b}\right\vert
h_{2}\psi_{d}\right)  \left.  =\right.  \sum_{1\leq k,l\leq2s}\left\{
h_{2lbmd}\left\vert \psi_{k}\right\rangle \left\langle \left.  \psi
_{l}\right\vert \psi_{c}\right\rangle \left\langle \psi_{a}\right\vert
+h_{2lamc}\left\vert \psi_{k}\right\rangle \left\langle \left.  \psi
_{l}\right\vert \psi_{d}\right\rangle \left\langle \psi_{b}\right\vert
\right\} \nonumber\\
&  \left.  =\right.  \sum_{1\leq k\leq2s}\left\{  h_{2kbcd}\left\vert \psi
_{k}\right\rangle \left\langle \psi_{a}\right\vert +h_{2kadc}\left\vert
\psi_{k}\right\rangle \left\langle \psi_{b}\right\vert \right\}  =\sum_{1\leq
m,n\leq2s}\left\{  h_{2mbcd}\delta_{an}+h_{2madc}\delta_{bn}\right\}
\left\vert \psi_{m}\right\rangle \left\langle \psi_{n}\right\vert .
\end{align}
This gives%
\begin{align}
&  \frac{\delta\mathcal{E}\left(  \mathbf{\psi}\right)  }{\delta\psi_{\mu
}^{\ast}}=\frac{1}{2}\sum_{1\leq m,n\leq2s}\left\{  \sum_{1\leq i,j,k,l}%
\left\{  \frac{1}{2N-1}\left\{  \overset{}{h_{1mk}\delta_{in}\left\vert
\psi_{m}\right\rangle \left\langle \psi_{n}\right\vert \delta_{jl}}%
-h_{1ml}\delta_{in}\left\vert \psi_{m}\right\rangle \left\langle \psi
_{n}\right\vert \delta_{jk}\right\}  \right.  \right. \\
&  \left.  \left.  \overset{}{\left\{  h_{2mjkl}\delta_{in}+h_{2milk}%
\delta_{jn}\right\}  \left\vert \psi_{m}\right\rangle \left\langle \psi
_{n}\right\vert -\left\{  h_{2mjlk}\delta_{in}+h_{2mikl}\delta_{jn}\right\}
}\right\}  D_{klij}^{2}\left\vert \psi_{m}\right\rangle \left\langle \psi
_{n}\right\vert \right\}  \left\vert \psi_{\mu}\right\rangle
\end{align}%
\begin{align}
&  \frac{\delta\mathcal{E}\left(  \mathbf{\psi}\right)  }{\delta\psi_{\mu
}^{\ast}}=\frac{1}{2}\sum_{1\leq m,n\leq2s}\left\{  \sum_{1\leq i,j,k,l}%
\left\{  \frac{1}{2N-1}\left\{  h_{1mk}\delta_{in}\delta_{jl}-h_{1ml}%
\delta_{in}\delta_{jk}\right\}  \right.  \right. \\
&  +\left.  \left.  \overset{}{\left\{  h_{2mjkl}\delta_{in}+h_{2milk}%
\delta_{jn}\right\}  -\left\{  h_{2mjlk}\delta_{in}+h_{2mikl}\delta
_{jn}\right\}  }\right\}  D_{klij}^{2}\left\vert \psi_{m}\right\rangle
\left\langle \psi_{n}\right\vert \right\}  \left\vert \psi_{\mu}\right\rangle
\end{align}%
\begin{align}
F\left(  \mathbf{\psi}\right)   &  =\frac{1}{2}\sum_{1\leq m,n\leq2s}\left\{
\sum_{1\leq i,j,k,l}\left\{  \frac{1}{\left(  2N-1\right)  }\left\{
h_{1mk}\delta_{in}\delta_{jl}-h_{1ml}\delta_{in}\delta_{jk}\right\}  \right.
\right. \\
&  +\left.  \left.  \overset{}{\left\{  h_{2mjkl}\delta_{in}+h_{2milk}%
\delta_{jn}\right\}  -\left\{  h_{2mjlk}\delta_{in}+h_{2mikl}\delta
_{jn}\right\}  }\right\}  D_{klij}^{2}\right\}  \left\vert \psi_{m}%
\right\rangle \left\langle \psi_{n}\right\vert
\end{align}%
\begin{align}
&  F\left(  \mathbf{\psi}\right)  =\frac{1}{2}\sum_{1\leq m,n\leq2s}\left\{
\frac{1}{\left(  2N-1\right)  }\sum_{1\leq j,k\leq2s}\left\{  h_{1mk}%
D_{kjnj}^{2}-h_{1mj}D_{jknj}^{2}\right\}  \right. \\
&  +\left.  \sum_{1\leq j,k,l\leq2}\left\{  \overset{}{\left\{  h_{2mjkl}%
D_{klnj}^{2}+h_{2mjlk}D_{kljn}^{2}\right\}  -\left\{  h_{2mjlk}D_{klnj}%
^{2}+h_{2mjkl}D_{kljn}^{2}\right\}  }\right\}  \right\}  \left\vert \psi
_{m}\right\rangle \left\langle \psi_{n}\right\vert
\end{align}%
\begin{equation}
F\left(  \mathbf{\psi}\right)  =\sum_{1\leq m,n\leq2s}\left\{  \frac
{1}{\left(  2N-1\right)  }\sum_{1\leq j,k\leq2s}h_{1mk}D_{kjnj}^{2}%
+\sum_{1\leq j,k,l\leq2s}\overset{}{\left\{  h_{2mjkl}D_{klnj}^{2}%
-h_{2mjlk}D_{klnj}^{2}\right\}  }\right\}  \left\vert \psi_{m}\right\rangle
\left\langle \psi_{n}\right\vert
\end{equation}%
\begin{equation}
F\left(  \mathbf{\psi}\right)  =\sum_{1\leq m,n\leq2s}\sum_{1\leq j,k,l\leq
2s}\left\{  \frac{1}{\left(  2N-1\right)  }h_{1mk}\delta_{lj}+\overset
{}{\left\{  h_{2mjkl}-h_{2mjlk}\right\}  }\right\}  D_{klnj}^{2}\left\vert
\psi_{m}\right\rangle \left\langle \psi_{n}\right\vert
\end{equation}%
\begin{align}
F\left(  \mathbf{\psi}\right)  ^{\dagger}  &  =\sum_{1\leq m,n\leq2s}%
\sum_{1\leq j,k,l\leq2s}\left\{  \frac{1}{\left(  2N-1\right)  }h_{1kn}%
\delta_{lj}+\overset{}{\left\{  h_{2klnj}-h_{2lknj}\right\}  }\right\}
D_{mjkl}^{2}\left\vert \psi_{m}\right\rangle \left\langle \psi_{n}\right\vert
\\
F\left(  \mathbf{\psi}\right)  _{mn}  &  =\sum_{1\leq j,k,l\leq2s}\left\{
\frac{1}{\left(  2N-1\right)  }h_{1mk}\delta_{lj}+\overset{}{\left\{
h_{2mjkl}-h_{2mjlk}\right\}  }\right\}  D_{klnj}^{2}\\
F\left(  \mathbf{\psi}\right)  _{nm}^{\ast}  &  =\sum_{1\leq j,k,l\leq
2s}\left\{  \frac{1}{\left(  2N-1\right)  }h_{1kn}\delta_{lj}+\overset
{}{\left\{  h_{2klnj}-h_{2lknj}\right\}  }\right\}  D_{mjkl}^{2}%
\end{align}%
\begin{align}
F_{2}\left(  \mathbf{\psi}\right)  _{mjkl}  &  =\frac{1}{\left(  2N-1\right)
}h_{1mk}\delta_{lj}+\overset{}{\left\{  h_{2mjkl}-h_{2mjlk}\right\}  }\\
F_{2}\left(  \mathbf{\psi}\right)  _{klmj}^{\ast}  &  =\frac{1}{\left(
2N-1\right)  }h_{1mk}\delta_{lj}+\overset{}{\left\{  h_{2mjkl}-h_{2mjlk}%
\right\}  =F_{2}\left(  \mathbf{\psi}\right)  _{mjkl}}%
\end{align}%
\begin{equation}
F\left(  \mathbf{\psi}\right)  =F\left(  \mathbf{\psi}\right)  ^{\dagger
}\Leftrightarrow L_{2}^{1}\left(  \left[  F_{2}\left(  \mathbf{\psi}\right)
,D^{2}\right]  \right)  =0
\end{equation}%
\begin{equation}
F\left(  \mathbf{\psi}\right)  =L_{2}^{1}\left(  F_{2}\left(  \mathbf{\psi
}\right)  D^{2}\right)  =\sum_{pq}L_{2}^{1}\left(  F_{2}\left(  \mathbf{\psi
}\right)  D^{2}\right)  _{pq}\left\vert \psi_{p}\right\rangle \left\langle
\psi_{q}\right\vert
\end{equation}%
\begin{equation}
F\left(  \mathbf{\psi}\right)  _{mm}=\sum_{1\leq j,k,l\leq2s}\left\{  \frac
{1}{\left(  2N-1\right)  }h_{1mk}\delta_{lj}+\overset{}{\left\{
h_{2mjkl}-h_{2mjlk}\right\}  }\right\}  D_{klmj}^{2}%
\end{equation}%
\begin{equation}
F\left(  \mathbf{\psi}\right)  _{mm}=\sum_{1\leq j\leq2s}\left\{  \frac
{1}{\left(  2N-1\right)  }h_{1mm}+\overset{}{\left\{  h_{2mjmj}-h_{2mjjm}%
\right\}  }\right\}  D_{mjmj}^{2}%
\end{equation}%
\begin{equation}
\sum_{m,n}c_{\nu m}^{\ast}F\left(  \mathbf{\psi}\right)  _{mn}c_{\lambda
n}=\sum_{m,n}\sum_{1\leq j,k,l\leq2s}\left\{  \frac{1}{\left(  2N-1\right)
}c_{\nu m}^{\ast}h_{1mk}\delta_{lj}+\overset{}{c_{\nu m}^{\ast}\left\{
h_{2mjkl}-h_{2mjlk}\right\}  }\right\}  D_{klnj}^{2}c_{\lambda n}%
\end{equation}


In second quantization%

\begin{align}
&  \left[  a_{r}^{\dagger}a_{s},a_{i}^{\dagger}a_{j}^{\dagger}a_{l}%
a_{k}\right]  =\left[  a_{r}^{\dagger}a_{s},a_{i}^{\dagger}\left(  \delta
_{jl}-a_{l}a_{j}^{\dagger}\right)  a_{k}\right]  =\left[  a_{r}^{\dagger}%
a_{s},a_{i}^{\dagger}a_{k}\right]  \delta_{jl}-\left[  a_{r}^{\dagger}%
a_{s},a_{i}^{\dagger}a_{l}a_{j}^{\dagger}a_{k}\right] \nonumber\\
&  =\left[  a_{r}^{\dagger}a_{s},a_{i}^{\dagger}a_{k}\right]  \delta
_{jl}-\left[  a_{r}^{\dagger}a_{s},a_{i}^{\dagger}a_{l}\right]  a_{j}%
^{\dagger}a_{k}-a_{i}^{\dagger}a_{l}\left[  a_{r}^{\dagger}a_{s}%
,a_{j}^{\dagger}a_{k}\right] \nonumber\\
&  =a_{r}^{\dagger}a_{k}\delta_{is}\delta_{jl}-a_{i}^{\dagger}a_{s}\delta
_{rk}\delta_{jl}-a_{r}^{\dagger}a_{l}a_{j}^{\dagger}a_{k}\delta_{is}%
+a_{i}^{\dagger}a_{s}a_{j}^{\dagger}a_{k}\delta_{rl}-a_{i}^{\dagger}a_{l}%
a_{r}^{\dagger}a_{k}\delta_{js}+a_{i}^{\dagger}a_{l}a_{j}^{\dagger}a_{s}%
\delta_{rk}%
\end{align}%
\begin{align}
&  i\left\langle \left.  \Psi\right\vert \left[  H,a_{r}^{\dagger}%
a_{s}\right]  \Psi\right\rangle =\frac{i}{4}\sum_{1\leq i,j,k,l\leq2s}%
K_{ijkl}\left\langle \left.  \Psi\right\vert \left[  a_{i}^{\dagger}%
a_{j}^{\dagger}a_{l}a_{k},a_{r}^{\dagger}a_{s}\right]  \Psi\right\rangle \\
&  =\frac{-i}{4}\sum_{1\leq i,j,k,l\leq2s}K_{ijkl}\left\langle \left.
\Psi\right\vert \left[  a_{r}^{\dagger}a_{s},a_{i}^{\dagger}a_{j}^{\dagger
}a_{l}a_{k}\right]  \Psi\right\rangle =\frac{-i}{4}\sum_{1\leq i,j,k,l\leq
2s}K_{ijkl}\times\\
&  \left\langle \left.  \Psi\right\vert \left\{  a_{r}^{\dagger}a_{k}%
\delta_{is}\delta_{jl}-a_{i}^{\dagger}a_{s}\delta_{rk}\delta_{jl}%
-a_{r}^{\dagger}a_{l}a_{j}^{\dagger}a_{k}\delta_{is}+a_{i}^{\dagger}a_{s}%
a_{j}^{\dagger}a_{k}\delta_{rl}-a_{i}^{\dagger}a_{l}a_{r}^{\dagger}a_{k}%
\delta_{js}+a_{i}^{\dagger}a_{l}a_{j}^{\dagger}a_{s}\delta_{rk}\right\}
\Psi\right\rangle \\
&  =\frac{-i}{4}\sum_{1\leq i,j,k,l\leq2s}K_{ijkl}\left\langle {}\right.
\left.  \Psi\right\vert \left\{  a_{r}^{\dagger}a_{k}\delta_{is}\delta
_{jl}-a_{r}^{\dagger}a_{k}\delta_{lj}\delta_{is}-a_{i}^{\dagger}a_{s}%
\delta_{rk}\delta_{jl}+a_{i}^{\dagger}a_{s}\delta_{rk}\delta_{lj}\right. \\
&  +a_{i}^{\dagger}a_{k}\delta_{rl}\delta_{sj}-a_{i}^{\dagger}a_{k}\delta
_{js}\delta_{rl}+\left.  a_{r}^{\dagger}a_{j}^{\dagger}a_{l}a_{k}\delta
_{is}+a_{i}^{\dagger}a_{r}^{\dagger}a_{l}a_{k}\delta_{js}-a_{i}^{\dagger}%
a_{j}^{\dagger}a_{s}a_{k}\delta_{rl}-a_{i}^{\dagger}a_{j}^{\dagger}a_{l}%
a_{s}\delta_{rk}\Psi\right\rangle \\
&  \frac{-i}{4}\sum_{1\leq j,k,l\leq2s}K_{ijkl}\left\langle \Psi|K_{sjkl}%
a_{r}^{\dagger}a_{j}^{\dagger}a_{l}a_{k}+K_{jskl}a_{j}^{\dagger}a_{r}%
^{\dagger}a_{l}a_{k}-K_{kjlr}a_{k}^{\dagger}a_{j}^{\dagger}a_{s}a_{l}%
-K_{kjrl}a_{k}^{\dagger}a_{j}^{\dagger}a_{l}a_{s}\Psi\right\rangle \\
&  =\frac{-i}{2}\sum_{1\leq j,k,l\leq2s}\left\langle \left.  \Psi\right\vert
\left\{  K_{sjkl}a_{r}^{\dagger}a_{j}^{\dagger}a_{l}a_{k}-K_{kjlr}%
a_{k}^{\dagger}a_{j}^{\dagger}a_{s}a_{l}\right\}  \Psi\right\rangle \\
&  =\frac{-i}{2}\sum_{1\leq j,k,l\leq2s}\left\{  K_{sjkl}D_{lkrj}^{2}%
-D_{lskj}^{2}K_{kjlr}\right\}
\end{align}%
\begin{equation}
D_{klij}^{2}=\operatorname{Tr}\left\{  a_{i}^{\dagger}a_{j}^{\dagger}%
a_{k}a_{l}D^{2N}\right\}
\end{equation}%
\begin{align}
&  =\frac{-i}{2}\sum_{1\leq j,k,l\leq2s}\left\{  K_{sjkl}D_{kljr}^{2}%
-D_{lskj}^{2}K_{kjlr}\right\}  =\frac{i}{2}\sum_{1\leq j,k,l\leq2s}\left\{
K_{sjkl}D_{klrj}^{2}-D_{slkj}^{2}K_{kjrl}\right\} \\
&  =\frac{i}{2}\sum_{j}\left\{  \left(  KD^{2}\right)  _{sjrj}-\left(
D^{2}K\right)  _{sjrj}\right\}  =\frac{i}{2}L_{2}^{1}\left(  \left[
K,D^{2}\right]  \right)  _{sr}%
\end{align}
Generalized Brillouin Condition for Optimized AGP States\label{BrillCon}

As discussed in \cite{weiner2004}\cite{Kurtz}\cite{Goscinski} one can define
operators $\left\{  \left.  q_{kj}^{\dagger}\right\vert -s\leq j\neq k\leq
s\right\}  $ by
\begin{equation}
q_{jk}^{\dagger}\left(  \mathbf{\eta,\varphi}\right)  =c_{k}\left(
\mathbf{\eta}\right)  a_{j}^{\dagger}a_{k}-\sigma\left(  j\right)
\sigma\left(  k\right)  c_{j}\left(  \mathbf{\eta}\right)  a_{-k}^{\dagger
}a_{-j}\text{ for }n_{j}\neq n_{k},
\end{equation}
where
\begin{equation}
\sigma\left(  k\right)  =sign\left(  k\right)  ,
\end{equation}
the operators $\left\{  a_{j}^{\dagger}\right\}  $ refer to basis formed by
the optimal CGSO's and which we have reindexed according to
\begin{equation}
\left\vert \varphi_{j}\right\rangle \rightarrow\left\vert \varphi
_{j-s}\right\rangle ;\quad1\leq j\leq2s.
\end{equation}
The importance of these operators lies in the observation that their adjoints
$\left\{  \left.  q_{kj}\right\vert -s\leq j\neq k\leq s\right\}  $ have the
special annihilation property%
\begin{equation}
q_{jk}\left\vert g\left(  \mathbf{\eta},\mathbf{\varphi}\right)
^{N}\right\rangle =0;\quad-s\leq j\neq k\leq s.
\end{equation}
One can show \cite{Weiner&Goscinski83} that
\begin{equation}
\left(
\begin{array}
[c]{c}%
q_{jk}^{\dagger}\\
q_{j-k}^{\dagger}\\
q_{-jk}^{\dagger}\\
q_{-j-k}^{\dagger}\\
q_{jk}\\
q_{j-k}\\
q_{-jk}\\
q_{-j-k}%
\end{array}
\right)  =\left(
\begin{array}
[c]{cccccccc}%
c_{k} & 0 & \cdots & \cdots & \cdots & \cdots & 0 & -c_{j}\\
0 & c_{k} & 0 & \cdots & \cdots & 0 & c_{j} & 0\\
\vdots & 0 & c_{k} & 0 & 0 & c_{j} & 0 & \vdots\\
\vdots & \vdots & 0 & c_{k} & -c_{j} & 0 & \vdots & \vdots\\
\vdots & \vdots & 0 & -c_{j} & c_{k} & 0 & \vdots & \vdots\\
\vdots & 0 & c_{j} & 0 & 0 & c_{k} & 0 & \vdots\\
0 & c_{j} & 0 & \cdots & \cdots & 0 & c_{k} & 0\\
-c_{j} & 0 & \cdots & \cdots & \cdots & \cdots & 0 & c_{k}%
\end{array}
\right)  \left(
\begin{array}
[c]{c}%
a_{j}^{\dagger}a_{k}\\
a_{j}^{\dagger}a_{-k}\\
a_{-j}^{\dagger}a_{k}\\
a_{-j}^{\dagger}a_{-k}\\
a_{k}^{\dagger}a_{j}\\
a_{-k}^{\dagger}a_{j}\\
a_{k}^{\dagger}a_{-j}\\
a_{-k}^{\dagger}a_{-j}%
\end{array}
\right)
\end{equation}
and that the transformation matrix is non singular. leading to the inverse
transformation
\begin{equation}
a_{j}^{\dagger}a_{k}=\sum_{l=-j,j;m=-k,k}\left\{  \alpha_{jklm}q_{lm}%
^{\dagger}+\beta_{jklm}q_{lm}\right\}  .
\end{equation}
The stationarity condition eq. $\left(  \ref{s3}\right)  $ can thus be
expressed as
\begin{equation}
\left\langle g\left(  \mathbf{\eta},\mathbf{\varphi}\right)  ^{N}\right.
\left\vert \left[  H,\sum_{l=-j,j;m=-k,k}\left\{  \alpha_{jklm}q_{lm}%
^{\dagger}+\beta_{jklm}q_{lm}\right\}  \right]  g\left(  \mathbf{\eta
},\mathbf{\varphi}\right)  ^{N}\right\rangle =0
\end{equation}
%

\begin{equation}
\left\langle g\left(  \mathbf{\eta},\mathbf{\varphi}\right)  ^{N}\right.
\left\vert H\sum_{l=-j,j;m=-k,k}\left\{  \alpha_{jklm}q_{lm}^{\dagger}%
+\beta_{jklm}q_{lm}\right\}  -H\sum_{l=-j,j;m=-k,k}\left\{  \alpha
_{jklm}q_{lm}^{\dagger}+\beta_{jklm}q_{lm}\right\}  g\left(  \mathbf{\eta
},\mathbf{\varphi}\right)  ^{N}\right\rangle =0
\end{equation}
%

\begin{equation}
\left\langle g\left(  \mathbf{\eta},\mathbf{\varphi}\right)  ^{N}\right.
\left\vert H\sum_{l=-j,j;m=-k,k}\alpha_{jklm}q_{lm}^{\dagger}g\left(
\mathbf{\eta},\mathbf{\varphi}\right)  ^{N}\right\rangle -\left\langle
g\left(  \mathbf{\eta},\mathbf{\varphi}\right)  ^{N}\right.  \left\vert
\sum_{l=-j,j;m=-k,k}\beta_{jklm}q_{lm}Hg\left(  \mathbf{\eta},\mathbf{\varphi
}\right)  ^{N}\right\rangle
\end{equation}%
\begin{equation}
a_{j}^{\dagger}a_{k}=\sum_{l=-j,j;m=-k,k}\left\{  \alpha_{jklm}u_{lm}%
+\beta_{jklm}v_{lm}\right\}
\end{equation}%
\begin{equation}
q_{kj}^{\dagger}=c_{j}a_{k}^{\dagger}a_{j}-\sigma\left(  j\right)
\sigma\left(  k\right)  c_{k}a_{-j}^{\dagger}a_{-k}%
\end{equation}%
\begin{equation}
q_{kj}^{\dagger}=c_{j}a_{k}^{\dagger}a_{j}-\sigma\left(  j\right)
\sigma\left(  k\right)  c_{k}a_{-j}^{\dagger}a_{-k}%
\end{equation}%
\begin{equation}
\mathcal{E}\left(  \mathbf{\psi}\right)  =\mathcal{E}\left(  \mathbf{z}%
_{+}\right)  =\frac{\left\langle c_{+1}\left(  \mathbf{z}_{+}\right)
c_{0}\left(  \mathbf{\lambda}\right)  g^{N}\right.  \left\vert Hc_{+1}\left(
\mathbf{z}_{+}\right)  c_{0}\left(  \mathbf{\lambda}\right)  g^{N}%
\right\rangle }{\left\langle c_{+1}\left(  \mathbf{z}_{+}\right)  c_{0}\left(
\mathbf{\lambda}\right)  g^{N}\right.  \left\vert c_{+1}\left(  \mathbf{z}%
_{+}\right)  c_{0}\left(  \mathbf{\lambda}\right)  g^{N}\right\rangle }%
\end{equation}%
\begin{align}
\mathcal{H}\left(  \mathbf{z}_{+}\right)   &  =\left\langle \exp\left\{
\sum_{kj}\mathbf{z}_{+kj}q_{kj}^{\dagger}\right\}  g_{0}^{N}\right.
\left\vert H\exp\left\{  \sum_{kj}\mathbf{z}_{+kj}q_{kj}^{\dagger}\right\}
g_{0}^{N}\right\rangle \\
&  \simeq\left\langle \left\{  1+\sum_{kj}\mathbf{z}_{+kj}q_{kj}^{\dagger
}\right\}  g_{0}^{N}\right.  \left\vert H\left\{  1+\sum_{kj}\mathbf{z}%
_{+kj}q_{kj}^{\dagger}\right\}  g_{0}^{N}\right\rangle \\
&  \simeq\left\langle g_{0}^{N}\right.  \left\vert Hg_{0}^{N}\right\rangle
+\sum_{kj}\left\{  \mathbf{z}_{+kj}^{\ast}\left\langle q_{kj}^{\dagger}%
g_{0}^{N}\right.  \left\vert Hg_{0}^{N}\right\rangle +\mathbf{z}%
_{+kj}\left\langle g_{0}^{N}\right.  \left\vert Hq_{kj}^{\dagger}g_{0}%
^{N}\right\rangle \right\} \\
&  =\mathcal{H}\left(  \mathbf{0}\right)  +\sum_{kj}\left\{  \mathbf{z}%
_{+kj}^{\ast}\left\langle g_{0}^{N}\right.  \left\vert q_{kj}Hg_{0}%
^{N}\right\rangle +\mathbf{z}_{+kj}\left\langle g_{0}^{N}\right.  \left\vert
Hq_{kj}^{\dagger}g_{0}^{N}\right\rangle \right\} \\
&  =\mathcal{H}\left(  \mathbf{0}\right)  +\sum_{kj}\left\{  \mathbf{z}%
_{+kj}^{\ast}\left\langle g_{0}^{N}\right.  \left\vert \left[  q_{kj}%
,H\right]  g_{0}^{N}\right\rangle +\mathbf{z}_{+kj}\left\langle g_{0}%
^{N}\right.  \left\vert \left[  H,q_{kj}^{\dagger}\right]  g_{0}%
^{N}\right\rangle \right\}
\end{align}%
\begin{equation}
\mathcal{E}\left(  \mathbf{\lambda}\right)  =\left\langle \exp\left\{
i\sum_{kj}\left\{  \mathbf{\lambda}_{kj}q_{kj}^{\dagger}+\mathbf{\lambda}%
_{kj}^{\ast}q_{kj}\right\}  \right\}  g_{0}^{N}\right.  \left\vert
H\exp\left\{  i\sum_{kj}\left\{  \mathbf{\lambda}_{kj}q_{kj}^{\dagger
}+\mathbf{\lambda}_{kj}^{\ast}q_{kj}\right\}  \right\}  g_{0}^{N}%
\right\rangle
\end{equation}%
\begin{align}
\delta^{1}\mathcal{E}\left(  \mathbf{0}\right)   &  =-\left\langle g_{0}%
^{N}\right.  \left\vert \left\{  i\sum_{kj}\left\{  \mathbf{\lambda}%
_{kj}^{\ast}q_{kj}+\mathbf{\lambda}_{kj}q_{kj}^{\dagger}\right\}  \right\}
Hg_{0}^{N}\right\rangle \\
&  \left\langle g_{0}^{N}\right.  \left\vert H\left\{  i\sum_{kj}\left\{
\mathbf{\lambda}_{kj}q_{kj}^{\dagger}+\mathbf{\lambda}_{kj}^{\ast}%
q_{kj}\right\}  \right\}  g_{0}^{N}\right\rangle \\
&  =i\sum_{kj}\left\{  \left\langle g_{0}^{N}\right.  \left\vert H\left\{
\mathbf{\lambda}_{kj}q_{kj}^{\dagger}+\mathbf{\lambda}_{kj}^{\ast}%
q_{kj}\right\}  g_{0}^{N}\right\rangle -\left\langle g_{0}^{N}\right.
\left\vert \left\{  \mathbf{\lambda}_{kj}^{\ast}q_{kj}+\mathbf{\lambda}%
_{kj}q_{kj}^{\dagger}\right\}  Hg_{0}^{N}\right\rangle \right\} \\
&  =i\sum_{kj}\left\{  \mathbf{\lambda}_{kj}\left\langle g_{0}^{N}\right.
\left\vert Hq_{kj}^{\dagger}g_{0}^{N}\right\rangle -\mathbf{\lambda}%
_{kj}^{\ast}\left\langle g_{0}^{N}\right.  \left\vert q_{kj}Hg_{0}%
^{N}\right\rangle \right\}  =0\forall\mathbf{\lambda}\\
&  =i\sum_{kj}\left\{  \left\langle g_{0}^{N}\right.  \left\vert
\mathbf{\lambda}_{kj}\left[  H,q_{kj}^{\dagger}\right]  +\mathbf{\lambda}%
_{kj}^{\ast}\left[  H,q_{kj}\right]  g_{0}^{N}\right\rangle \right\}
=0\forall\mathbf{\lambda}%
\end{align}%
\begin{align}
&  \left\langle g_{0}^{N}\right.  \left\vert Hq_{kj}^{\dagger}g_{0}%
^{N}\right\rangle -\left\langle g_{0}^{N}\right.  \left\vert q_{kj}Hg_{0}%
^{N}\right\rangle =0\\
&  \left\langle g_{0}^{N}\right.  \left\vert Hq_{kj}^{\dagger}g_{0}%
^{N}\right\rangle +\left\langle g_{0}^{N}\right.  \left\vert q_{kj}Hg_{0}%
^{N}\right\rangle =0\\
&  \Rightarrow\left\langle g_{0}^{N}\right.  \left\vert Hq_{kj}^{\dagger}%
g_{0}^{N}\right\rangle =0
\end{align}%
\begin{align}
a_{j}^{\dagger}a_{k}  &  =\sum_{l=-j,j;m=-k,k}\left\{  \alpha_{jklm}%
q_{lm}^{\dagger}+\beta_{jklm}q_{lm}\right\}  \Rightarrow\\
\left\langle g_{0}^{N}\right.  \left\vert Ha_{j}^{\dagger}a_{k}g_{0}%
^{N}\right\rangle  &  =\left\langle g_{0}^{N}\right.  \left\vert
H\sum_{l=-j,j;m=-k,k}\left\{  \alpha_{jklm}q_{lm}^{\dagger}+\beta_{jklm}%
q_{lm}\right\}  g_{0}^{N}\right\rangle \\
&  =\left\langle g_{0}^{N}\right.  \left\vert H\sum_{l=-j,j;m=-k,k}%
\alpha_{jklm}q_{lm}^{\dagger}g_{0}^{N}\right\rangle =0
\end{align}
Stationarity condition $\mathbf{z}_{+}$
\begin{equation}
\sum_{kj}\left\{  \mathbf{z}_{+kj}^{\ast}\left\langle g_{0}^{N}\right.
\left\vert q_{kj}Hg_{0}^{N}\right\rangle +\mathbf{z}_{+kj}\left\langle
g_{0}^{N}\right.  \left\vert Hq_{kj}^{\dagger}g_{0}^{N}\right\rangle \right\}
=0
\end{equation}
must be true for all
\begin{align}
&  \left\langle g_{0}^{N}\right.  \left\vert \left[  H,q_{kj}^{\dagger
}\right]  g_{0}^{N}\right\rangle =0\\
&  =\left\langle g_{0}^{N}\right.  \left\vert \left[  H,c_{j}a_{k}^{\dagger
}a_{j}-\sigma\left(  j\right)  \sigma\left(  k\right)  c_{k}a_{-j}^{\dagger
}a_{-k}\right]  g_{0}^{N}\right\rangle =0\\
&  =c_{j}\left\langle g_{0}^{N}\right.  \left\vert \left[  H,a_{k}^{\dagger
}a_{j}\right]  \left.  g_{0}^{N}\right\rangle -\sigma\left(  j\right)
\sigma\left(  k\right)  c_{k}\left\langle g_{0}^{N}\right.  \left[
H,a_{-j}^{\dagger}a_{-k}\right]  g_{0}^{N}\right\rangle =0\\
&  \Leftrightarrow c_{j}L_{2}^{1}\left(  \left[  K,D^{2}\right]  \right)
_{jk}-\sigma\left(  j\right)  \sigma\left(  k\right)  c_{k}L_{2}^{1}\left(
\left[  K,D^{2}\right]  \right)  _{-k-j}\,=0\forall\quad\text{unpaired }k,j
\end{align}%
\begin{align}
&  c_{j}L_{2}^{1}\left(  KD^{2}\right)  _{jk}-\sigma\left(  j\right)
\sigma\left(  k\right)  c_{k}L_{2}^{1}\left(  KD^{2}\right)  _{-k-j}\\
&  c_{j}f_{jk}-\sigma\left(  j\right)  \sigma\left(  k\right)  c_{k}f_{-k-j}%
\end{align}%
\begin{equation}
F\left(  \mathbf{\psi}\right)  =L_{2}^{1}\left(  F_{2}\left(  \mathbf{\psi
}\right)  D^{2}\right)
\end{equation}%
\begin{equation}
\exp\left\{  \sum_{kj}\mathbf{z}_{+kj}q_{kj}^{\dagger}\right\}  \left\vert
g_{0}\right\rangle =e^{\Lambda}\left\vert g_{0}\right\rangle =\sum_{1\leq
j\leq s}n_{j}^{\frac{1}{2}}e^{\Lambda}\left\vert \chi_{j}\chi_{j+s}%
\right\rangle =\sum_{1\leq j\leq s}n_{j}^{\frac{1}{2}}\left\vert \chi
_{j}\left(  \mathbf{z}_{+}\right)  \chi_{j+s}\left(  \mathbf{z}_{+}\right)
\right\rangle
\end{equation}%
\begin{equation}
\exp\left\{  \sum_{kj}\mathbf{z}_{+kj}q_{kj}^{\dagger}\right\}  \left\vert
g_{0}\right\rangle =e^{\Lambda}\left\vert g_{0}\right\rangle =\sum_{1\leq
j\leq s}\xi_{j}e^{\Lambda}\left\vert \chi_{j}\chi_{j+s}\right\rangle
=\sum_{1\leq j\leq s}\xi_{j}\left\vert \chi_{j}\left(  \mathbf{z}_{+}\right)
\chi_{j+s}\left(  \mathbf{z}_{+}\right)  \right\rangle
\end{equation}%
\begin{equation}
c_{0}\left(  \mathbf{\lambda}\right)  =\exp\left\{  \sum
\limits_{_{\substack{1\leq j\leq s\\1\leq k\leq4}}}\lambda_{jk}n_{jk}\right\}
=c_{0-1}\left(  \lambda_{j2}\right)  c_{00}\left(  \mathbf{\lambda}%
_{j0}\right)  c_{01}\left(  \lambda_{j4}\right)
\end{equation}%
\begin{align}
\exp\left\{  \sum_{1\leq j\leq s}\lambda_{j1}n_{j1}\right\}  \left\vert
g_{0}\right\rangle  &  =\sum_{1\leq j\leq s}\xi_{j}\left(  \lambda
_{j1}\right)  \left\vert \chi_{j}\chi_{j+s}\right\rangle \\
\lambda_{j1}  &  =c_{j}e^{i\theta_{j}}\\
\xi_{j}\left(  \lambda_{j1}\right)   &  =\xi_{j}\left(  c_{j},\theta
_{j}\right)  =c_{j}e^{i\theta_{j}}c_{0j}%
\end{align}%
\begin{align}
n_{j1}  &  =a_{j}^{\dagger}a_{j}+a_{-j}^{\dagger}a_{-j}\nonumber\\
n_{j2}  &  =a_{j}^{\dagger}a_{-j}+a_{-j}^{\dagger}a_{j}\\
n_{j3}  &  =a_{j}^{\dagger}a_{j}-a_{-j}^{\dagger}a_{-j}\\
n_{j4}  &  =i\left(  a_{j}^{\dagger}a_{-j}-a_{-j}^{\dagger}a_{j}\right)
\end{align}%
\begin{align}
c_{0-1}\left(  \lambda_{j2}\right)   &  =e^{\lambda_{j2}n_{j2}}\\
c_{01}\left(  \lambda_{j2}\right)   &  =e^{\lambda_{j2}n_{j4}}\\
c_{0-1}\left(  \lambda_{j2}\right)   &  =e^{\lambda_{j1}n_{j1}+\lambda
_{j3}n_{j3}}%
\end{align}%
\begin{equation}
c_{0-1}\left(  \lambda_{j2}\right)  c_{01}\left(  \lambda_{j4}\right)
e^{\lambda_{j3}n_{j3}}\left\vert g_{0}\right\rangle =\left\vert g_{0}%
\right\rangle
\end{equation}
For the paired operators
\begin{align}
\left\langle g_{0}^{N}\right.  \left\vert \left[  H,n_{j2}\right]  g_{0}%
^{N}\right\rangle  &  =0\\
\left\langle g_{0}^{N}\right.  \left\vert \left[  H,n_{j3}\right]  g_{0}%
^{N}\right\rangle  &  =0\\
\left\langle g_{0}^{N}\right.  \left\vert \left[  H,n_{j4}\right]  g_{0}%
^{N}\right\rangle  &  =0
\end{align}
while
\begin{equation}
\left\langle g_{0}^{N}\right.  \left\vert \left[  H,n_{j1}\right]  g_{0}%
^{N}\right\rangle =0
\end{equation}
gives the stationary phase condition and
\begin{equation}
\left\langle g_{0}^{N}\right.  \left\vert Hn_{1j}g_{0}^{N}\right\rangle
=\mathcal{E}\left\langle g_{0}^{N}\right.  \left\vert n_{j1}g_{0}%
^{N}\right\rangle
\end{equation}
gives the stationary occupation number condition. This comes from
\begin{equation}
\frac{1}{S_{N}}\left\{  \left\langle g_{0}^{N}\right.  \left\vert \left(
n_{1j}H+Hn_{1j}\right)  g_{0}^{N}\right\rangle -2\mathcal{E}\left\langle
g_{0}^{N}\right.  \left\vert n_{1j}g_{0}^{N}\right\rangle \right\}  =0.
\end{equation}%
\begin{align}
&  \left.  \left\langle \exp\left\{  \sum_{j}\mathbf{z}_{+j}n_{1j}\right\}
g_{0}^{N}\right.  \left\vert H\exp\left\{  \sum_{j}\mathbf{z}_{+j}%
n_{1j}\right\}  g_{0}^{N}\right\rangle \right/  S_{N}\\
&  =\left.  \left\langle \exp\left\{  \sum_{j}\left(  \alpha_{j}+i\beta
_{j}\right)  n_{1j}\right\}  g_{0}^{N}\right.  \left\vert H\exp\left\{
\sum_{j}\left(  \alpha_{j}+i\beta_{j}\right)  n_{1j}\right\}  g_{0}%
^{N}\right\rangle \right/  S_{N}\\
&  \Rightarrow\\
\delta^{1}\mathcal{E}  &  =\sum_{j}\left\langle g_{0}^{N}\right.  \left\vert
\left\{  \alpha_{j}\left(  n_{1j}H+Hn_{1j}\right)  +i\beta_{j}\left[
H,n_{1j}\right]  -2\alpha_{j}\mathcal{E}n_{1j}\right\}  g_{0}^{N}%
\right\rangle
\end{align}
This gives the condition for stationary phase as
\begin{equation}
\left\langle g_{0}^{N}\right.  \left\vert \left[  H,n_{1j}\right]  g_{0}%
^{N}\right\rangle =0
\end{equation}
and for stationary occupation as
\begin{align}
&  \left\langle g_{0}^{N}\right.  \left\vert \left\{  n_{1j}H+Hn_{1j}%
-2\alpha_{j}\mathcal{E}n_{1j}\right\}  g_{0}^{N}\right\rangle =0\\
&  \Rightarrow\left\langle g_{0}^{N}\right.  \left\vert \left\{
n_{1j}H+Hn_{1j}\right\}  g_{0}^{N}\right\rangle =2\mathcal{E}\left\langle
g_{0}^{N}\right.  \left\vert n_{1j}g_{0}^{N}\right\rangle \\
&  \Rightarrow\left\langle g_{0}^{N}\right.  \left\vert \left\{
n_{1j}H\right\}  g_{0}^{N}\right\rangle =\mathcal{E}\left\langle g_{0}%
^{N}\right.  \left\vert n_{1j}g_{0}^{N}\right\rangle
\end{align}
Ionization energies associated with the GSO's $\left\{  \left\vert \psi
_{k}\right\rangle \right\}  $ are given by
\begin{equation}
\frac{\left\langle g^{N}\right.  \left\vert Hg^{N}\right\rangle }{\left\langle
g^{N}\right.  \left\vert g^{N}\right\rangle }-\frac{\left\langle a_{k}%
g^{N}\right.  \left\vert Ha_{k}g^{N}\right\rangle }{\left\langle a_{k}%
g^{N}\right.  \left\vert a_{k}g^{N}\right\rangle }\equiv\frac{\left\langle
g^{N}\right.  \left\vert Hg^{N}\right\rangle }{\left\langle g^{N}\right.
\left\vert g^{N}\right\rangle }-\frac{\left\langle g^{N}\right.  \left\vert
a_{k}^{\dagger}Ha_{k}g^{N}\right\rangle }{\left\langle g^{N}\right.
\left\vert a_{k}^{\dagger}a_{k}g^{N}\right\rangle }%
\end{equation}
when the AGP\ is optimized this gives
\begin{equation}
\mathcal{E-}\frac{\left\langle g^{N}\right.  \left\vert a_{k}^{\dagger}%
Ha_{k}g^{N}\right\rangle }{\left\langle g^{N}\right.  \left\vert
a_{k}^{\dagger}a_{k}g^{N}\right\rangle }\equiv\frac{\left\langle g^{N}\right.
\left\vert H\left(  a_{k}^{\dagger}a_{k}+a_{-k}^{\dagger}a_{-k}\right)
g^{N}\right\rangle }{n_{k}}-\frac{\left\langle g^{N}\right.  \left\vert
a_{k}^{\dagger}Ha_{k}g^{N}\right\rangle }{n_{k}}=
\end{equation}


\section{Derivatives of the Energy Functional\label{variational}}

The $2N$-AGP\ energy functional in terms of the CGSO's $\left\{  \left.
\psi_{j}\right\vert 1\leq j\leq2s\right\}  $, the geminal occupation numbers
$\left\{  \left.  n_{j}\right\vert 1\leq j\leq s\right\}  $ and the phases
$\left\{  \left.  \theta_{j}\right\vert 1\leq j\leq s\right\}  $ can be
expressed as the quotient%
\begin{align}
&  \mathcal{E}\left(  \mathbf{n},\mathbf{\psi},\mathbf{\theta}\right)
=\frac{\mathcal{H}\left(  \mathbf{n},\mathbf{\psi},\mathbf{\theta}\right)
}{S_{N}\left(  \mathbf{n}\right)  }\nonumber\\
&  =\frac{1}{S_{N}\left(  \mathbf{n}\right)  }\left\{  \sum_{1\leq j\leq
s}n_{j}\left\{  \left\langle \psi_{j}\left\vert h_{1}\psi_{j}\right.
\right\rangle +\left\langle \psi_{j+s}\left\vert h_{1}\psi_{j+s}\right.
\right\rangle +\psi_{j}\psi_{j+s}|\left\vert \psi_{j}\psi_{j+s}\right\rangle
\right\}  \right. \nonumber\\
&  +\sum_{1\leq j,k\leq s}n_{j}^{\frac{1}{2}}n_{k}^{\frac{1}{2}}e^{i\left(
\theta_{j}-\theta_{k}\right)  }S_{N-1}\left(  \jmath k\right)  \psi_{j}%
\psi_{j+s}|\left\vert \psi_{k}\psi_{k+s}\right\rangle \left.  +\frac{1}{2}%
\sum_{1\leq j,k\leq s}n_{j}n_{k}S_{N-2}\left(  \jmath k\right)  \left\langle
\psi_{j}\psi_{k}|||\psi_{j}\psi_{k}\right\rangle \right\}  .
\end{align}
The first order variation $\delta^{1}\mathcal{E}\left(  \mathbf{n}%
,\mathbf{\psi},\mathbf{\theta}\right)  $ of the energy is
\begin{equation}
\delta^{1}\mathcal{E}\left(  \mathbf{n},\mathbf{\psi},\mathbf{\theta}\right)
=\left(
\begin{array}
[c]{cccc}%
\partial_{\mathbf{n}}^{1}\mathcal{E}\left(  \mathbf{n},\mathbf{\psi
},\mathbf{\theta}\right)  & \partial_{\mathbf{\psi}^{\ast}}^{1}\mathcal{E}%
\left(  \mathbf{n},\mathbf{\psi},\mathbf{\theta}\right)  & \partial
_{\mathbf{\psi}}^{1}\mathcal{E}\left(  \mathbf{n},\mathbf{\psi},\mathbf{\theta
}\right)  & \partial_{\mathbf{\theta}}^{1}\mathcal{E}\left(  \mathbf{n}%
,\mathbf{\psi},\mathbf{\theta}\right)
\end{array}
\right)  \left(
\begin{array}
[c]{c}%
\delta\mathbf{n}\\
\delta\mathbf{\psi}^{\ast}\\
\delta\mathbf{\psi}\\
\delta\mathbf{\theta}%
\end{array}
\right)  ,
\end{equation}
where
\begin{subequations}
\label{Fderiv}%
\begin{align}
\partial_{\mathbf{n}}^{1}\mathcal{E}\left(  \mathbf{n},\mathbf{\psi
},\mathbf{\theta}\right)   &  =\frac{1}{S_{N}\left(  \mathbf{n}\right)
}\left\{  \partial_{\mathbf{n}}^{1}\mathcal{H}\left(  \mathbf{n},\mathbf{\psi
},\mathbf{\theta}\right)  -\mathcal{E}\left(  \mathbf{n},\mathbf{\psi
},\mathbf{\theta}\right)  \partial_{\mathbf{n}}^{1}S_{N}\left(  \mathbf{n}%
\right)  \right\} \\
\partial_{\mathbf{\psi}^{\ast}}^{1}\mathcal{E}\left(  \mathbf{n},\mathbf{\psi
},\mathbf{\theta}\right)   &  =\frac{1}{S_{N}\left(  \mathbf{n}\right)
}\left\{  \partial_{\mathbf{\psi}^{\ast}}^{1}\mathcal{H}\left(  \mathbf{n}%
,\mathbf{\psi},\mathbf{\theta}\right)  -\mathcal{E}\left(  \mathbf{n}%
,\mathbf{\psi},\mathbf{\theta}\right)  \partial_{\mathbf{\psi}^{\ast}}%
^{1}S_{N}\left(  \mathbf{n}\right)  \right\}  =\left[  \partial_{\mathbf{\psi
}}^{1}\mathcal{E}\left(  \mathbf{n},\mathbf{\psi},\mathbf{\theta}\right)
\right]  ^{\ast}\\
\partial_{\mathbf{\theta}}^{1}\mathcal{E}\left(  \mathbf{n},\mathbf{\psi
},\mathbf{\theta}\right)   &  =\frac{1}{S_{N}\left(  \mathbf{n}\right)
}\left\{  \partial_{\mathbf{\theta}}^{1}\mathcal{H}\left(  \mathbf{n}%
,\mathbf{\psi},\mathbf{\theta}\right)  -\mathcal{E}\left(  \mathbf{n}%
,\mathbf{\psi},\mathbf{\theta}\right)  \partial_{\mathbf{\theta}}^{1}%
S_{N}\left(  \mathbf{n}\right)  \right\}
\end{align}
and
\end{subequations}
\begin{subequations}
\begin{align}
\partial_{\mathbf{\psi}^{\ast}}^{1}\mathcal{E}\left(  \mathbf{n},\mathbf{\psi
},\mathbf{\theta}\right)  \cdot\delta\mathbf{\psi}^{\ast}  &  =\sum_{1\leq
j\leq2s}\partial_{\psi_{j}^{\ast}}^{1}\mathcal{E}\left(  \mathbf{n}%
,\mathbf{\psi},\mathbf{\theta}\right)  \left(  \delta\psi_{j}^{\ast}\right)
=\sum_{1\leq j\leq2s}\int\frac{\delta\mathcal{E}\left(  \mathbf{n}%
,\mathbf{\psi},\mathbf{\theta}\right)  }{\delta\psi_{j}^{\ast}\left(
\mathbf{r,}\zeta\right)  }\delta\psi_{j}^{\ast}\left(  \mathbf{r,}%
\zeta\right)  d\mathbf{r}d\zeta\label{funcder}\\
\partial_{\mathbf{n}}^{1}\mathcal{E}\left(  \mathbf{n},\mathbf{\psi
},\mathbf{\theta}\right)  \cdot\delta\mathbf{n}  &  =\sum_{1\leq j\leq s}%
\frac{\partial\mathcal{E}\left(  \mathbf{n},\mathbf{\psi},\mathbf{\theta
}\right)  }{\partial n_{j}}\delta n_{j}\\
\partial_{\mathbf{\theta}}^{1}\mathcal{E}\left(  \mathbf{n},\mathbf{\psi
},\mathbf{\theta}\right)  \cdot\delta\mathbf{\theta}  &  =\sum_{1\leq j\leq
s}\frac{\partial\mathcal{E}\left(  \mathbf{n},\mathbf{\psi},\mathbf{\theta
}\right)  }{\partial\theta_{j}}\delta\theta_{j}.
\end{align}
If the variations $\left\{  \delta\mathbf{\psi}^{\ast},\delta\mathbf{\psi
}\right\}  $ maintain the orthonormality of the CGSO's i.e.
\end{subequations}
\begin{equation}
\left\langle \left.  \psi_{j}\right\vert \psi_{k}\right\rangle =\delta_{jk}%
\end{equation}
then
\begin{subequations}
\label{fderiv2}%
\begin{align}
\partial_{\mathbf{\psi}^{\ast}}^{1}\mathcal{E}\left(  \mathbf{n},\mathbf{\psi
},\mathbf{\theta}\right)   &  =\frac{1}{S_{N}\left(  \mathbf{n}\right)
}\partial_{\mathbf{\psi}^{\ast}}^{1}\mathcal{H}\left(  \mathbf{n}%
,\mathbf{\psi},\mathbf{\theta}\right) \\
\partial_{\mathbf{\theta}}^{1}\mathcal{E}\left(  \mathbf{n},\mathbf{\psi
},\mathbf{\theta}\right)   &  =\frac{1}{S_{N}\left(  \mathbf{n}\right)
}\partial_{\mathbf{\theta}}^{1}\mathcal{H}\left(  \mathbf{n},\mathbf{\psi
},\mathbf{\theta}\right)
\end{align}
as $\partial_{\mathbf{\psi}^{\ast}}^{1}S_{N}\left(  \mathbf{n}\right)
=\partial_{\mathbf{\psi}}^{1}S_{N}\left(  \mathbf{n}\right)  =\partial
_{\mathbf{\theta}}^{1}S_{N}\left(  \mathbf{n}\right)  =\mathbf{0}$.

Variation in terms of the variables $\left\{  \mathbf{n},\mathbf{\psi
},\mathbf{\theta}\right\}  $ leads to a constrained variational problem, which
can be avoided if one uses the group parameters $\left\{  \mathbf{\eta
},\mathbf{\lambda},\mathbf{\theta}\right\}  $ and the energy functional
\end{subequations}
\begin{equation}
\mathcal{E}\left(  \mathbf{\eta},\mathbf{\lambda},\mathbf{\theta}\right)
=\frac{\mathcal{H}\left(  \mathbf{\eta},\mathbf{\lambda},\mathbf{\theta
}\right)  }{S_{N}\left(  \mathbf{\eta}\right)  }=\frac{\left\langle \left.
g\left(  \mathbf{\eta},\mathbf{\lambda},\mathbf{\theta}\right)  ^{N}%
\right\vert Hg\left(  \mathbf{\eta},\mathbf{\lambda},\mathbf{\theta}\right)
^{N}\right\rangle }{S_{N}\left(  \mathbf{\eta}\right)  }%
\end{equation}
and considers the derivatives
\begin{subequations}
\label{GFderiv}%
\begin{align}
\partial_{\mathbf{\eta}}^{1}\mathcal{E}\left(  \mathbf{\eta},\mathbf{\lambda
},\mathbf{\theta}\right)   &  =\frac{1}{S_{N}\left(  \mathbf{\eta}\right)
}\left\{  \partial_{\mathbf{\eta}}^{1}\mathcal{H}\left(  \mathbf{\eta
},\mathbf{\lambda},\mathbf{\theta}\right)  -\mathcal{E}\left(  \mathbf{\eta
},\mathbf{\lambda},\mathbf{\theta}\right)  \partial_{\mathbf{\eta}}^{1}%
S_{N}\left(  \mathbf{\eta}\right)  \right\}  \\
\partial_{\mathbf{\lambda}}^{1}\mathcal{E}\left(  \mathbf{\eta}%
,\mathbf{\lambda},\mathbf{\theta}\right)   &  =\frac{1}{S_{N}\left(
\mathbf{\eta}\right)  }\partial_{\mathbf{\lambda}}^{1}\mathcal{H}\left(
\mathbf{\eta},\mathbf{\lambda},\mathbf{\theta}\right)  \\
\partial_{\mathbf{\theta}}^{1}\mathcal{E}\left(  \mathbf{\eta},\mathbf{\lambda
},\mathbf{\theta}\right)   &  =\frac{1}{S_{N}\left(  \mathbf{\eta}\right)
}\partial_{\mathbf{\theta}}^{1}\mathcal{H}\left(  \mathbf{\eta}%
,\mathbf{\lambda},\mathbf{\theta}\right)  .
\end{align}
The generators of the GSO and the phase variations produce the derivatives
\end{subequations}
\begin{equation}
\partial_{\mathbf{\lambda}_{jk}}^{1}\mathcal{E}\left(  \mathbf{\eta
},\mathbf{\lambda},\mathbf{\theta}\right)  =\frac{i}{S_{N}\left(
\mathbf{\eta}\right)  }\left\langle \left.  g\left(  \mathbf{\eta
},\mathbf{\lambda},\mathbf{\theta}\right)  ^{N}\right\vert \left[
H,a_{j}^{\dagger}a_{k}\right]  g\left(  \mathbf{\eta},\mathbf{\lambda
},\mathbf{\theta}\right)  ^{N}\right\rangle ;\quad1\leq j\neq k\pm
s\leq2s\label{Dlamda}%
\end{equation}
and
\begin{equation}
\partial_{\mathbf{\theta}_{j}}^{1}\mathcal{E}\left(  \mathbf{\eta
},\mathbf{\lambda},\mathbf{\theta}\right)  =\frac{i}{S_{N}\left(
\mathbf{\eta}\right)  }\left\langle \left.  g\left(  \mathbf{\eta
},\mathbf{\lambda},\mathbf{\theta}\right)  ^{N}\right\vert \left[  H,\left(
a_{j}^{\dagger}a_{j}+a_{j+s}^{\dagger}a_{j+s}\right)  \right]  g\left(
\mathbf{\eta},\mathbf{\lambda},\mathbf{\theta}\right)  ^{N}\right\rangle
;\quad1\leq j\leq s,\label{Dtheta}%
\end{equation}
while the generators of the occupation number variations give the derivatives
\begin{align}
\partial_{\mathbf{\eta}_{j}}^{1}\mathcal{E}\left(  \mathbf{\eta}%
,\mathbf{\lambda},\mathbf{\theta}\right)   &  =\frac{1}{S_{N}\left(
\mathbf{\eta}\right)  }\left\{  \left\langle \left.  g\left(  \mathbf{\eta
},\mathbf{\lambda},\mathbf{\theta}\right)  ^{N}\right\vert \left[  H,\left(
a_{j}^{\dagger}a_{j}+a_{j+s}^{\dagger}a_{j+s}\right)  \right]  _{+}g\left(
\mathbf{\eta},\mathbf{\lambda},\mathbf{\theta}\right)  ^{N}\right\rangle
\right.  \nonumber\\
&  \left.  -2\mathcal{E}\left(  \mathbf{\eta},\mathbf{\lambda},\mathbf{\theta
}\right)  \left\langle \left.  g\left(  \mathbf{\eta},\mathbf{\lambda
},\mathbf{\theta}\right)  ^{N}\right\vert \left(  a_{j}^{\dagger}a_{j}%
+a_{j+s}^{\dagger}a_{j+s}\right)  g\left(  \mathbf{\eta},\mathbf{\lambda
},\mathbf{\theta}\right)  ^{N}\right\rangle \right\}  ;\quad1\leq j\leq
s\label{Deta}%
\end{align}
Density Matrix Functional\label{DMF}

The AGP states that have the same FORDO are parameterized by a vector
$\mathbf{\theta}$ and have energies given by
\begin{equation}
\tilde{f}\left(  \mathbf{\theta}\right)  =K\left(  \mathbf{n},\mathbf{\psi
}\right)  +\frac{1}{S_{N}\left(  \mathbf{n}\right)  }\sum_{1\leq j,k\leq
s}n_{j}^{\frac{1}{2}}n_{k}^{\frac{1}{2}}e^{i\left(  \theta_{j}-\theta
_{k}\right)  }S_{N-1}\left(  \jmath k\right)  \left\langle \psi_{j}\psi
_{j+s}|h_{2}\psi_{k}\psi_{k+s}\right\rangle ,
\end{equation}
where the constant $K$ is given by eq. $\left(  \ref{thetaen}\right)  $ and
\begin{equation}
\mathbf{\theta}=\left(  0,\theta_{2},\cdots,\theta_{s}\right)  ;\quad
0<\theta_{j}\leq2\pi,\quad2\leq j\leq s.
\end{equation}
As $K\left(  \mathbf{n},\mathbf{\psi}\right)  $ is constant for fixed $\left(
\mathbf{n},\mathbf{\psi}\right)  $ we are interested in the minimum of
\begin{align}
f\left(  \mathbf{\theta}\right)   &  =\frac{1}{S_{N}\left(  \mathbf{n}\right)
}\sum_{1\leq j,k\leq s}n_{j}^{\frac{1}{2}}n_{k}^{\frac{1}{2}}e^{i\left(
\theta_{j}-\theta_{k}\right)  }S_{N-1}\left(  \jmath k\right)  \left\langle
\psi_{j}\psi_{j+s}\right\vert \left\vert \psi_{k}\psi_{k+s}\right\rangle
\nonumber\\
&  =\sum_{1\leq j,k\leq s}e^{-i\theta_{k}}M_{kj}e^{i\theta_{j}}=\left(
e^{i\mathbf{\theta}}\right)  ^{\dagger}\mathbb{M}e^{i\mathbf{\theta}},
\end{align}
where $e^{i\mathbf{\theta}}$ is the vector
\begin{align}
\left(  e^{i\mathbf{\theta}}\right)  _{j}  &  =e^{i\theta_{j}};\quad2\leq
j\leq s\nonumber\\
\left(  e^{i\mathbf{\theta}}\right)  _{1}  &  =1
\end{align}
and $\mathbb{M}$ the positive matrix given by
\begin{equation}
M_{kj}=\frac{1}{S_{N}\left(  \mathbf{n}\right)  }n_{j}^{\frac{1}{2}}%
n_{k}^{\frac{1}{2}}S_{N-1}\left(  \jmath k\right)  \left\langle \psi_{j}%
\psi_{j+s}\right\vert \left\vert \psi_{k}\psi_{k+s}\right\rangle .
\end{equation}
We can expand $f$ in a Taylors series about $\mathbf{\theta}_{0}$ to obtain
\begin{equation}
f\left(  \mathbf{\theta}_{0}+\delta\mathbf{\theta}\right)  =f\left(
\mathbf{\theta}_{0}\right)  +\partial_{\mathbf{\theta}}f\left(  \mathbf{\theta
}_{0}\right)  ^{t}\delta\mathbf{\theta}+\frac{1}{2}\delta\mathbf{\theta}%
^{t}\partial_{\mathbf{\theta}}^{2}f\left(  \mathbf{\theta}_{0}\right)
\delta\mathbf{\theta+\cdots,}%
\end{equation}
where
\begin{equation}
\partial_{\mathbf{\theta}}f\left(  \mathbf{\theta}_{0}\right)  _{l}%
=\frac{\partial f\left(  \mathbf{\theta}_{0}\right)  }{\partial\theta_{l}%
}=i\sum_{1\leq j\leq s}\left\{  e^{-i\theta_{0j}}M_{jl}e^{i\theta_{0l}%
}-e^{-i\theta_{0l}}M_{lj}e^{i\theta_{0j}}\right\}  =-2\sum_{1\leq j\leq
s}\operatorname{Im}\left\{  e^{-i\theta_{0j}}M_{jl}e^{i\theta_{0l}}\right\}
\end{equation}
and
\begin{align}
\partial_{\mathbf{\theta}}^{2}f\left(  \mathbf{\theta}_{0}\right)  _{lm}  &
=\frac{\partial^{2}f\left(  \mathbf{\theta}_{0}\right)  }{\partial\theta
_{l}\partial\theta_{m}}\nonumber\\
&  =2\left\{  \operatorname{Im}\left\{  \sum_{1\leq j\leq s}e^{-i\theta_{0j}%
}M_{jl}e^{i\theta_{0l}}\right\}  \delta_{lm}-\operatorname{Im}\left\{
e^{-i\theta_{0l}}M_{lm}e^{i\theta_{0m}}\right\}  \right\}  .
\end{align}
The stationary points $\mathbf{\theta}_{\#}$ of $f$ are given by
\begin{equation}
\partial_{\mathbf{\theta}}f\left(  \mathbf{\theta}_{\#}\right)  =\mathbf{0}
\label{gradf}%
\end{equation}
and these are local minimum if
\begin{equation}
\partial_{\mathbf{\theta}}^{2}f\left(  \mathbf{\theta}_{\#}\right)
\geq\mathbf{0.}%
\end{equation}
The eq. $\left(  \ref{gradf}\right)  $ can be solved numerically by any
standard optimization technique, though it would be more convenient if an
analytic solution to eq. $\left(  \ref{gradf}\right)  $ could be found.

\subsection{Stationarity Conditions}

A necessary condition for $\left(  \mathbf{n}_{\#}\mathbf{,\psi}_{\#}\right)
$ to be a solution of eq. $\left(  \ref{min}\right)  $ is that there exists
Lagrange multipliers $\left\{  \left.  \varepsilon_{jk}\right\vert 1\leq
j,k\leq2s\right\}  $ such that
\begin{equation}
\partial^{1}\mathcal{L}\left(  \mathbf{n}_{\#}\mathbf{,\psi}_{\#}\right)
=\mathbf{0,}%
\end{equation}
where
\begin{equation}
\left.  \partial^{1}\mathcal{L}\left(  \mathbf{n}_{\#}\mathbf{,\psi}%
_{\#}\right)  \right.  =\left(  \partial_{\mathbf{n}}^{1}\mathcal{L}\left(
\mathbf{n}_{\#}\mathbf{,\psi}_{\#}\right)  ,\partial_{\mathbf{\psi}}%
^{1}\mathcal{L}\left(  \mathbf{n}_{\#}\mathbf{,\psi}_{\#}\right)  \right)
\end{equation}
and
\begin{align}
\partial_{\mathbf{n}}^{1}\mathcal{\tilde{E}}\left(  \mathbf{n}_{\#}%
\mathbf{,\psi}_{\#}\right)  _{j}  &  =\frac{\partial\mathcal{\tilde{E}}\left(
\mathbf{n}_{\#}\mathbf{,\psi}_{\#}\right)  }{\partial n_{j}};\quad1\leq j\leq
s\nonumber\\
\partial_{\mathbf{\psi}}^{1}\mathcal{\tilde{E}}\left(  \mathbf{n}%
_{\#}\mathbf{,\psi}_{\#}\right)  _{j}  &  =\frac{\delta\mathcal{\tilde{E}%
}\left(  \mathbf{n}_{\#}\mathbf{,\psi}_{\#}\right)  }{\delta\psi_{j}^{\ast}%
};\quad1\leq j\leq s\nonumber\\
\partial_{\mathbf{\psi}}^{1}\mathcal{\tilde{E}}\left(  \mathbf{n}%
_{\#}\mathbf{,\psi}_{\#}\right)  _{j+s}  &  =\frac{\delta\mathcal{\tilde{E}%
}\left(  \mathbf{n}_{\#}\mathbf{,\psi}_{\#}\right)  }{\delta\psi_{j}}%
;\quad1\leq j\leq s
\end{align}
and that $\left(  \mathbf{n}_{\#}\mathbf{,\psi}_{\#}\right)  $ satisfy the
constraints eqs. $\left(  \ref{pos1}\right)  ,\left(  \ref{gemnorm}\right)  $
and $\left(  \ref{orthonorm}\right)  $.As is well known this condition only
characterizes constrained stationary points of $\mathcal{E}$ and the Hessian
matrix of second derivatives has to be examined in order to determine if the
point $\left(  \mathbf{n}_{\#}\mathbf{,\psi}_{\#}\right)  $ does in fact
minimize $\mathcal{E}.$

\section{Optimized Orbital Fock Operator\label{saddle}}

The Lagrangian that controls the constrained minimization of the energy,
$\mathcal{E}\left(  \mathbb{D}^{2}\mathbf{,\psi}\right)  $, wrt the GSO's
$\left\{  \left.  \left\vert \psi_{j}\right\rangle \right\vert 1\leq
j\leq2s\right\}  $ is
\begin{equation}
\mathcal{L}\left(  \mathbb{D}^{2}\mathbf{,\psi,\varepsilon}\right)
=\mathcal{E}\left(  \mathbb{D}^{2}\mathbf{,\psi}\right)  -\sum_{1\leq
j,k\leq2s}\varepsilon_{kj}\left(  \left\langle \left.  \psi_{j}\right\vert
\psi_{k}\right\rangle -\delta_{jk}\right)  .
\end{equation}
This can be reformulated as
\begin{equation}
\mathcal{L}\left(  \mathbb{D}^{2}\mathbf{,\psi,\varepsilon}\right)
=\mathcal{E}\left(  \mathbb{D}^{2}\mathbf{,\psi}\right)  -\operatorname{Tr}%
\left\{  \mathbf{\varepsilon c}\left(  \mathbf{\psi}^{\ast},\mathbf{\psi
}\right)  \right\}  ,
\end{equation}
where
\begin{align}
\mathbf{\varepsilon}  &  \mathbb{\leftrightarrow}\left\{  \varepsilon
_{kj}\right\}  ,\nonumber\\
\mathbf{c}\left(  \mathbf{\psi}^{\ast},\mathbf{\psi}\right)   &
\mathbb{\leftrightarrow}\left\{  c_{kj}\right\}
\end{align}
and
\begin{equation}
c_{kj}=\left\langle \left.  \psi_{j}\right\vert \psi_{k}\right\rangle
-\delta_{jk}.
\end{equation}
The Lagrange multiplier term can be formulated as a sesquilinear functional of
$\left(  \mathbf{\psi}^{\ast},\mathbf{\psi}\right)  $ and
\begin{equation}
f_{\mathbf{\varepsilon}}\left(  \mathbf{\psi}^{\ast},\mathbf{\psi}\right)
=\operatorname{Tr}\left\{  \mathbf{\varepsilon c}\left(  \mathbf{\psi}^{\ast
},\mathbf{\psi}\right)  \right\}
\end{equation}
a linear functional of $\mathbf{\varepsilon}$%
\begin{equation}
f_{\mathbf{c}\left(  \mathbf{\psi}^{\ast},\mathbf{\psi}\right)  }\left(
\mathbf{\varepsilon}\right)  =\operatorname{Tr}\left\{  \mathbf{\varepsilon
c}\left(  \mathbf{\psi}^{\ast},\mathbf{\psi}\right)  \right\}  .
\end{equation}
The Kuhn-Tucker conditions \cite{KuhnTuck} show that if $\left(  \mathbf{\psi
}_{\#},\mathbf{\varepsilon}_{\#}\right)  $ describes a constrained minimum of
$\mathcal{E}$ then $\left(  \mathbf{\psi}_{\#},\mathbf{\varepsilon}%
_{\#}\right)  $ is a saddle point of $\mathcal{L}\left(  \mathbb{D}%
^{2}\mathbf{,\psi,\varepsilon}\right)  ,$ such that
\begin{equation}
\mathcal{L}\left(  \mathbb{D}^{2}\mathbf{,\mathbf{\psi}_{\#}%
,\mathbf{\varepsilon}_{\#}+\delta\varepsilon}\right)  \leq\mathcal{L}\left(
\mathbb{D}^{2}\mathbf{,\psi}_{\#},\mathbf{\varepsilon}_{\#}\right)
\leq\mathcal{L}\left(  \mathbb{D}^{2}\mathbf{,\mathbf{\psi}_{\#}%
+\delta\mathbf{\psi},\mathbf{\varepsilon}_{\#}}\right)
\end{equation}
for any changes $\left\{  \mathbf{\delta\mathbf{\psi}}^{\ast}%
\mathbf{\mathbf{,\delta\mathbf{\psi,}}\delta\varepsilon}\right\}  .$ The
Taylors series expansion in $\left\{  \varepsilon_{kj}\right\}  $ gives
\begin{equation}
\mathcal{L}\left(  \mathbb{D}^{2}\mathbf{,\psi}_{\#},\mathbf{\varepsilon}%
_{\#}+\delta\mathbf{\varepsilon}\right)  =\mathcal{E}\left(  \mathbb{D}%
^{2}\mathbf{,\psi}_{\#}\right)  -\operatorname{Tr}\left\{  \left(
\mathbf{\varepsilon}_{\#}+\delta\mathbf{\varepsilon}\right)  \mathbf{c}\left(
\mathbf{\psi}_{\#}^{\ast},\mathbf{\psi}_{\#}\right)  \right\}  =\mathcal{L}%
\left(  \mathbb{D}^{2}\mathbf{,\psi}_{\#},\mathbf{\varepsilon}_{\#}\right)
-\operatorname{Tr}\left\{  \delta\mathbf{\varepsilon c}\left(  \mathbf{\psi
}_{\#}^{\ast},\mathbf{\psi}_{\#}\right)  \right\}  .
\end{equation}
The Taylors series expansion in $\left(  \mathbf{\psi}^{\ast},\mathbf{\psi
}\right)  $ gives
\begin{align}
&  \mathcal{L}\left(  \mathbb{D}^{2}\mathbf{,\psi}_{\#}+\delta\mathbf{\psi
},\mathbf{\varepsilon}_{\#}\right)  =\mathcal{E}\left(  \mathbb{D}%
^{2}\mathbf{,\psi}_{\#}+\delta\mathbf{\psi}\right)  -\operatorname{Tr}\left\{
\mathbf{\varepsilon}_{\#}\mathbf{c}\left(  \mathbf{\psi}_{\#}^{\ast
},\mathbf{\psi}_{\#}\right)  \right\} \\
&  =\mathcal{E}\left(  \mathbb{D}^{2}\mathbf{,\psi}_{\#}+\delta\mathbf{\psi
}\right)  -\operatorname{Tr}\left\{  \mathbf{\varepsilon}_{\#}\mathbf{c}%
\left(  \mathbf{\psi}_{\#}^{\ast}+\delta\mathbf{\psi}^{\ast},\mathbf{\psi
}_{\#}+\delta\mathbf{\psi}\right)  \right\} \\
&  =\mathcal{L}\left(  \mathbb{D}^{2}\mathbf{,\psi}_{\#},\mathbf{\varepsilon
}_{\#}\right)  +\partial^{2}\mathcal{E}\left(  \mathbb{D}^{2}\mathbf{,\psi
}_{\#}\right)  \left(  \delta\mathbf{\psi}^{\ast}\mathbf{,}\delta\mathbf{\psi
}\right)  -\operatorname{Tr}\left\{  \mathbf{\varepsilon}_{\#}\mathbf{c}%
\left(  \delta\mathbf{\psi}^{\ast},\delta\mathbf{\psi}\right)  \right\}
\geq\mathcal{L}\left(  \mathbb{D}^{2}\mathbf{,\psi}_{\#},\mathbf{\varepsilon
}_{\#}\right) \\
&  \Rightarrow\partial^{2}\mathcal{E}\left(  \mathbb{D}^{2}\mathbf{,\psi}%
_{\#}\right)  \left(  \delta\mathbf{\psi}^{\ast}\mathbf{,}\delta\mathbf{\psi
}\right)  \geq\operatorname{Tr}\left\{  \mathbf{\varepsilon}_{\#}%
\mathbf{c}\left(  \delta\mathbf{\psi}^{\ast},\delta\mathbf{\psi}\right)
\right\}  =\sum\mathbf{\varepsilon}_{\#jk}\left\langle \left.  \delta\psi
_{k}\right\vert \delta\psi_{j}\right\rangle =\delta\mathbf{\psi}^{\dagger
}\mathbf{\varepsilon}_{\#}\delta\mathbf{\psi}\\
&  \Rightarrow\partial^{2}\mathcal{E}\left(  \mathbb{D}^{2}\mathbf{,\psi}%
_{\#}\right)  \geq\mathbf{\varepsilon}_{\#}\Rightarrow\delta\mathbf{\psi
}\partial^{2}\mathcal{E}\left(  \mathbb{D}^{2}\mathbf{,\psi}_{\#}\right)
\delta\mathbf{\psi\geq}\delta\mathbf{\psi}^{\dagger}\mathbf{\varepsilon}%
_{\#}\delta\mathbf{\psi\quad\forall}\delta\mathbf{\psi}%
\end{align}
At a minimum of $\mathcal{E}$ one has $\partial^{2}\mathcal{E}\left(
\mathbb{D}^{2}\mathbf{,\psi}_{\#}\right)  \geq\mathbf{0}$

We can express the derivatives of $2N-1$ electron states $\left\{  a_{p}%
D^{2N}a_{p}^{\dagger}\right\}  $ in terms of a generalized Fock operator%
\begin{equation}
\frac{\delta\mathcal{E}_{p}\left(  D^{2N},\mathbf{\varphi}\right)  }%
{\delta\varphi_{j}^{\ast}}=F\left(  \mathbb{D}^{2}\left(  p\right)
,\mathbf{\varphi}\right)  \left\vert \varphi_{j}\right\rangle ,
\end{equation}
where
\begin{equation}
D^{2}\left(  p\right)  =\frac{1}{\operatorname{Tr}\left\{  a_{p}^{\dagger
}a_{p}D^{2N}\right\}  }L_{N}^{2}\left(  a_{p}D^{2N}a_{p}^{\dagger}\right)  .
\end{equation}
The derivatives wrt occupation number formally give
\begin{equation}
\frac{\partial\mathcal{E}\left(  \mathbb{D}^{2},\mathbf{\varphi}\right)
}{\partial N_{j}}=\frac{1}{2\left(  N-1\right)  }\left\{  \left\{
\mathcal{E}_{p}\left(  D^{2N},\mathbf{\varphi}\right)  +\left\langle \left.
\varphi_{p}\right\vert h_{1}\varphi_{p}\right\rangle \right\}  +\frac
{\partial\mathcal{E}_{p}\left(  D^{2N},\mathbf{\varphi}\right)  }{\partial
N_{j}}\right\}  =0.
\end{equation}
In the case of AGP states occupation number variations can be implemented by
the group $U_{occnum}$ giving
\begin{align}
&  \frac{\partial\mathcal{E}\left(  \mathbb{D}^{2},\mathbf{\varphi}\right)
}{\partial\eta_{j}}=\frac{1}{2\left(  N-1\right)  }\sum_{1\leq p\leq
2s}\left\{  \frac{\partial N_{p}\left(  \mathbf{\eta}\right)  }{\partial
\eta_{j}}\left\{  \mathcal{E}_{p}\left(  k\left(  \mathbf{\eta}\right)
D^{2N}k\left(  \mathbf{\eta}\right)  ,\mathbf{\varphi}\right)  +\left\langle
\left.  \varphi_{p}\right\vert h_{1}\varphi_{p}\right\rangle \right\}  \right.
\\
&  \left.  +\right.  \left.  N_{p}\left(  \mathbf{\eta}\right)  \frac
{\partial\mathcal{E}_{p}\left(  k\left(  \mathbf{\eta}\right)  D^{2N}k\left(
\mathbf{\eta}\right)  ,\mathbf{\varphi}\right)  }{\partial\eta_{j}}\right\}
=0
\end{align}%
\begin{align}
\partial_{\mathbf{\eta}_{j}}^{1}\mathcal{E}\left(  \mathbf{\eta}%
,\mathbf{\lambda},\mathbf{\theta}\right)   &  =\frac{1}{S_{N}\left(
\mathbf{\eta}\right)  }\left\{  \left\langle \left.  g\left(  \mathbf{\eta
},\mathbf{\lambda},\mathbf{\theta}\right)  ^{N}\right\vert \left[  H,\left(
a_{j}^{\dagger}a_{j}+a_{j+s}^{\dagger}a_{j+s}\right)  \right]  _{+}g\left(
\mathbf{\eta},\mathbf{\lambda},\mathbf{\theta}\right)  ^{N}\right\rangle
\right. \nonumber\\
&  -\left.  2\mathcal{E}\left(  \mathbf{\eta},\mathbf{\lambda},\mathbf{\theta
}\right)  \left\langle \left.  g\left(  \mathbf{\eta},\mathbf{\lambda
},\mathbf{\theta}\right)  ^{N}\right\vert \left(  a_{j}^{\dagger}a_{j}%
+a_{j+s}^{\dagger}a_{j+s}\right)  g\left(  \mathbf{\eta},\mathbf{\lambda
},\mathbf{\theta}\right)  ^{N}\right\rangle \right\}  =0;\quad1\leq j\leq s,
\end{align}%
\begin{align}
&  \frac{\partial N_{p}\left(  \mathbf{\eta}\right)  }{\partial\eta_{j}}%
=\frac{\partial}{\partial\eta_{j}}\left\{  \frac{1}{S_{N}\left(  \mathbf{\eta
}\right)  }\left\langle \left.  g\left(  \mathbf{\eta},\mathbf{\lambda
},\mathbf{\theta}\right)  ^{N}\right\vert a_{p}^{\dagger}a_{p}g\left(
\mathbf{\eta},\mathbf{\lambda},\mathbf{\theta}\right)  ^{N}\right\rangle
\right\} \\
&  =\frac{1}{S_{N}\left(  \mathbf{\eta}\right)  }\left\{  \left\langle \left.
g\left(  \mathbf{\eta},\mathbf{\lambda},\mathbf{\theta}\right)  ^{N}%
\right\vert \left\{  a_{j}^{\dagger}a_{j}a_{p}^{\dagger}a_{p}+a_{p}^{\dagger
}a_{p}a_{j}^{\dagger}a_{j}\right\}  g\left(  \mathbf{\eta},\mathbf{\lambda
},\mathbf{\theta}\right)  ^{N}\right\rangle \right. \\
&  \left.  -2N_{p}\left(  \mathbf{\eta}\right)  \left\langle \left.  g\left(
\mathbf{\eta},\mathbf{\lambda},\mathbf{\theta}\right)  ^{N}\right\vert
\left\{  a_{j}^{\dagger}a_{j}\right\}  g\left(  \mathbf{\eta},\mathbf{\lambda
},\mathbf{\theta}\right)  ^{N}\right\rangle \right\}
\end{align}%
\begin{align}
\frac{\partial\mathcal{E}_{p}\left(  k\left(  \mathbf{0}\right)
D^{2N}k\left(  \mathbf{0}\right)  ,\mathbf{\varphi}\right)  }{\partial\eta
_{j}}  &  =\frac{1}{N_{p}}\left\{  \left\langle \left.  g\left(  \mathbf{\eta
},\mathbf{\varphi}\right)  ^{N}\right\vert \left[  a_{p}^{\dagger}%
Ha_{p},\left(  a_{j}^{\dagger}a_{j}+a_{j+s}^{\dagger}a_{j+s}\right)  \right]
_{+}g\left(  \mathbf{\eta},\mathbf{\varphi}\right)  ^{N}\right\rangle \right.
\\
&  -\mathcal{E}_{p}\left(  \left(  \mathbf{\eta},\mathbf{\varphi}\right)
\right)  \left\langle \left.  g\left(  \mathbf{\eta},\mathbf{\varphi}\right)
^{N}\right\vert \left[  a_{p}^{\dagger}a_{p},\left(  a_{j}^{\dagger}%
a_{j}+a_{j+s}^{\dagger}a_{j+s}\right)  \right]  _{+}g\left(  \mathbf{\eta
},\mathbf{\varphi}\right)  ^{N}\right\rangle
\end{align}
In terms of CGSO's $\mathbf{\psi}$ the AGP-SORDM can be expressed solely in
terms of the geminal occupation numbers i.e. $\mathbb{D}^{2}=\mathbb{D}%
^{2}\left(  \mathbf{n}\right)  $ and if these occupation numbers are also
optimized and non degenerate the Fock operator $F\left(  \mathbf{n}%
,\mathbf{\psi}\right)  $ can be diagonalized to obtain eigenvalues that are
again weighted ionization potentials associated with eigenvectors $\left\{
\left\vert \varsigma_{p}\right\rangle \right\}  $, but these eigenvectors are
not NGSO's nor CGSO's of the optimized AGP state.


\end{document}